\documentclass[11pt,a4paper]{article}
\usepackage[T1]{fontenc}
\usepackage[utf8]{inputenc}
\usepackage{lmodern,microtype}
\usepackage[margin=.85in]{geometry}
\usepackage{amsmath,amssymb,amsthm,mathtools,aliascnt}
\usepackage{booktabs,tabularx,array,enumitem,needspace,float}
\usepackage[section]{placeins}
\usepackage{xcolor,tikz,mdframed}
\usetikzlibrary{arrows.meta,positioning,calc,decorations.pathreplacing}
\usepackage[hidelinks]{hyperref}
\usepackage[nameinlink,noabbrev]{cleveref}
\usepackage{fancyhdr}
\hypersetup{pdftitle={Sync HotStuff RR, Part II: Sequencer Rotation},pdfsubject={Sequencer rotation with regional and global finality}}
\setlength{\headheight}{14pt}
\pagestyle{fancy}\fancyhf{}
\lhead{\small Sync HotStuff RR, Part II}\rhead{\small Sequencer rotation and finality}
\cfoot{\thepage}
\setlength{\parindent}{0pt}\setlength{\parskip}{4pt}
\setlength{\emergencystretch}{2em}
\renewcommand{\topfraction}{.92}
\renewcommand{\textfraction}{.06}
\renewcommand{\floatpagefraction}{.85}
\makeatletter
\setlength{\@fptop}{0pt}\setlength{\@fpbot}{0pt plus 1fil}
\makeatother
\setlist{nosep,leftmargin=*}
\newcolumntype{Y}{>{\raggedright\arraybackslash}X}
\definecolor{rrRed}{HTML}{A51C30}
\definecolor{rrBlue}{HTML}{2563EB}
\definecolor{rrTeal}{HTML}{0F766E}
\newcommand{\RR}[1]{{\color{rrRed}#1}}
\newcommand{\QC}{\mathcal C}
\newcommand{\FC}{\mathcal F}
\newcommand{\Hash}{\mathsf H}
\newcommand{\rank}{\operatorname{rank}}
\newcommand{\Del}{\operatorname{Del}}
\newcommand{\Twait}{T_{\rm wait}}
\newcommand{\TB}{\mathit{TB}}
\newcommand{\tagg}[1]{\ensuremath{\mathsf{#1}}}
\newcommand{\defterm}[1]{\textcolor{rrBlue}{\textit{#1}}}
\newtheorem{definition}{Definition}[section]
\newaliascnt{assumption}{definition}
\newtheorem{assumption}[assumption]{Assumption}
\aliascntresetthe{assumption}
\crefname{assumption}{assumption}{assumptions}
\Crefname{assumption}{Assumption}{Assumptions}
\newcounter{protocol}
% Protocol boxes float to the top of a page, so long boxes leave no gaps.
\newenvironment{protocol}[2]{\begin{figure}[htbp]\refstepcounter{protocol}\label{#1}
\begin{mdframed}[linewidth=.6pt,skipabove=0pt,skipbelow=0pt,
innertopmargin=7pt,innerbottommargin=7pt]\small
\textbf{Protocol \theprotocol. #2}\par\smallskip}{\end{mdframed}\end{figure}}
\crefname{protocol}{Protocol}{Protocols}
\Crefname{protocol}{Protocol}{Protocols}
\title{Sync HotStuff RR, Part II: Sequencer Rotation\\\large Multi-Region Rotation with Regional and Global Finality}
\author{}\date{13 October 2026}
\begin{document}
\maketitle
\begin{abstract}
Sync HotStuff RR lets several regions take turns as the
sequencer of one chain. One turn of one region is a \defterm{session}.
Inside a session, the region runs a stake-weighted version of the
Part~I protocol~\cite{part1} until it commits a fixed quota of DATA blocks;
the next live region then continues from the last block (regular
handoff). If a region fails, an honest multisignature committee posts a
Force on L1. The Force cuts the session's external history at a chosen
height and hands the turn to another region (forced handoff). A block is
\defterm{regionally final} when it has a valid commit receipt from the
region and session that the rotation history assigns to it. It becomes
\defterm{globally final} at an observer after a further waiting period,
which starts when the observer relays the complete commit evidence to
the committee. The period covers the committee's monitoring, its
decision window, and the L1 finalization of any Force, so a valid Force
never removes a globally final block. A Force passes three stages:
Submitted, L1-finalized, and L1-activated. Finalization stops the
affected blocks from becoming global; activation switches the external
branch, and the old regional log stays unchanged. We specify the model, six
protocols, and a conditional safety argument. A full liveness and
complexity analysis is outside the scope of this document.
\end{abstract}

\section{Overview}
\label{sec:overview}
Part~I~\cite{part1} adapts standard-synchronous Sync
HotStuff~\cite{sync} to one fixed committee with scheduled leader
rotation. This document runs that protocol inside each region and lets
several regions take turns producing blocks. \defterm{Session} $r$ is one turn,
served by region $i_r$.

\paragraph{What this document adds to Part~I.}
Part~I orders blocks for a single committee. Its leaders rotate every
$K$ blocks, and its static Byzantine adversary corrupts at most $f$ of
$n=2f+1$ replicas. Part~II keeps that engine and builds a sequencer
service around it. The new features are:
\begin{itemize}
\item \textbf{Sessions and a three-level view.} Regions take turns. One
turn is a session with a fixed quota $\TB_r$ of DATA blocks. Part~I's
view $(g,t)$ becomes $(r,g,t)$, which adds the session index $r$.
\item \textbf{Regular handoff without L1.} When the session's last quota
block $T_r$ has a regional receipt, the next live region starts from
$T_r$. No L1 transaction and no global wait are needed.
\item \textbf{Forced handoff through L1.} If a region fails, an honest
multisignature committee posts a Force on L1. The Force cuts the
session's external history at a height $h_*$ and selects another
region. A Force passes three stages: Submitted, L1-finalized, and
L1-activated.
\item \textbf{Two finality levels.} A block is regionally final when a
receipt from its assigned session proves its commitment. It becomes
globally final after a waiting period $\Twait$. The committee's
protected prefix guarantees that no valid Force removes a globally final
block.
\item \textbf{Session IDs and branches.} A Force can replace a regular
boundary, so sessions form branches. A session is named by a hash of the
boundary that opened it, not by its index.
\end{itemize}
Inside a region, Protocols~1 and~2 change Part~I in a few places.
Certificates are stake-weighted. Terms stop at the session quota. A
\defterm{DRAIN} block commits the quota's last block after a recovery.
Replicas send \defterm{FINAL} messages that form regional receipts, and a
new \defterm{Completion} step ends the session. 
This FINAL-Completion stage ensures that old region ends the regional session only after the external observice regionally finalize the last-session block.
  
Recovery entry also waits for
the next generation's authority. \Cref{sec:protocols} lists these changes,
and \RR{red text} marks them in the protocol boxes. Protocols~3--6, which
handle rotation, Forces, and observers, are new in full and are not
marked.

\Cref{tab:roles} lists the participants. A \defterm{node} is any
participant that follows the external chain; every replica is a node.
Each node is also an \defterm{observer}: it keeps its own global prefix
(\cref{sec:external}). We write ``observer $u$'' when a rule depends on
$u$'s local state.

\begin{table}[htbp]\centering\small
\caption{Participants and their responsibilities.}
\label{tab:roles}
\begin{tabularx}{\linewidth}{@{}>{\raggedright\arraybackslash}p{.15\linewidth}>{\raggedright\arraybackslash}p{.25\linewidth}Y@{}}
\toprule
Role & Who & Responsibility\\
\midrule
\defterm{L1} & The settlement chain & Public source of authority:
voters, leaders, and quotas of each session; status updates and Forces
(\cref{sec:model}).\\
\defterm{Replica} & A member of a region's voting set $V_i$ &
Runs the regional engine of each session of its region
(\cref{alg:steady,alg:recovery}); enters a new session when its region
is selected (\cref{alg:bootstrap}).\\
\defterm{Committee} & One honest multisignature service &
Monitors regional output, keeps a protected prefix, and posts status
updates and Forces (\cref{sec:force}).\\
\defterm{Node} (\defterm{observer}) & Every replica and every external verifier &
Follows the current branch, processes L1 records, and maintains
regional and global finality (\cref{alg:rotation}, steps 2--3 of
\cref{alg:force}, and \cref{alg:external}).\\
\defterm{Backbone} & The network between regions & Eventually
delivers blocks, receipts, and witnesses. It never chooses the next
region; L1 records and regional certificates do.\\
\bottomrule
\end{tabularx}
\end{table}

A session ends by a \defterm{regular handoff}, after its last quota block
$T_r$ commits, or by a \defterm{forced handoff}, when an L1 force
instruction cuts its external history at a DATA height $h_*$.
\Cref{sec:examples} gives both paths. Either ending fixes a
\defterm{boundary}: the retained anchor block, the next region, and the L1
height that selected it (\cref{sec:witness}). A Force can replace a
regular boundary, even after the regular successor has started; the
sessions after the replaced boundary are then discarded. Sessions
therefore form \defterm{branches}. A session is identified by its
\defterm{session ID} $\sigma_r$, the hash of the boundary that opened it,
and not by its index $r$, which a Force may reuse. A node's
\defterm{current branch} is the one selected by regular boundaries and
activated Forces.

Finality has two levels (\cref{sec:external}). A block is
\defterm{regionally final} when the rotation history assigns it to a region and
session, and a valid receipt from that session proves its commitment. A
receipt alone does not prove that the producing region was authorized.
A block becomes \defterm{globally final} at an observer after a waiting period
that starts when the observer relays the complete commit evidence to the
committee, never at the block's older L1 reference. Regular handoffs use
regional finality only; they never wait for global finality.

A \defterm{view} is $(r,g,t)$: \defterm{session} $r$, \defterm{recovery generation} $g$, and \defterm{leader
term} $t$. Views on one branch are ordered lexicographically. Each session starts at
$(r,0,0)$. \Cref{tab:views} lists the four ways to change view. The
first two happen inside region $i_r$, as in Part~I. The members of the
next region take the last two.

\begin{table}[htbp]\centering\small
\caption{The four view changes.}
\label{tab:views}
\begin{tabularx}{\linewidth}{@{}>{\raggedright\arraybackslash}p{.19\linewidth}>{\raggedright\arraybackslash}p{.12\linewidth}Y>{\raggedright\arraybackslash}p{.23\linewidth}@{}}
\toprule
Change & New view & Trigger & Commit timers\\
\midrule
\defterm{Normal rotation} & $(r,g,t+1)$ & A terminal QC of term $t$, for DATA below $H_r$ & Keep running.\\
\defterm{Recovery} & $(r,g+1,0)$ & A quit proof for generation $g$, then a $\Delta$ wait, and $\tagg{AUTH}_{r,g+1}$ & Canceled in generation $g$.\\
\defterm{Regular handoff} & $(r+1,0,0)$ & Ready on a REGULAR witness, which holds a receipt for $T_r$ (\cref{sec:ready}) & The old engine cancels its own timers in Completion (\cref{alg:steady}).\\
\defterm{Forced handoff} & $(r+1,0,0)$ & Ready on a FORCE witness, which holds an activated force & The old engine is unchanged.\\
\bottomrule
\end{tabularx}
\end{table}
A QC for DATA at $H_r$, or for a DRAIN, starts no new term. A status
update without a force changes no view. Recovery keeps the session's
quota. A Force may replace a regular boundary. The new
session has a different session ID, even if it reuses the numeric
index.

\section{Model and assumptions}
\label{sec:model}

\subsection{L1 control state}
\label{sec:l1}
At each finalized L1 height, nodes read the same authority records and
status-update list (\cref{tab:l1}).
Finalized L1 block headers also give their authenticated timestamps.
The multisignature committee posts status updates. A regular handoff
needs no new L1 transaction.

\begin{table}[htbp]\centering\small
\caption{L1 records used by the protocol.}
\label{tab:l1}
\begin{tabularx}{\linewidth}{@{}p{.27\linewidth}Y@{}}
\toprule
Field & Meaning\\
\midrule
\multicolumn{2}{@{}l}{\textbf{\defterm{Authority}} $\tagg{AUTH}_{r,g}$}\\
\defterm{Stake configuration} & The members $V_i$ of each region, their public
keys and voting weights, the Byzantine bound $F_i$ used as
the certificate threshold, and protocol parameters such as $K$ and
$\Delta_i$.\\
\defterm{Leader schedule} & The leaders within each region. Once a boundary
chooses $i_r$, this gives $L(r,g,t)\in V_{i_r}$.\\
\defterm{Regional quota} $\TB_r$ & The number of new DATA blocks that session $r$
may add. Recovery does not reset it.\\
\midrule
\multicolumn{2}{@{}l}{\textbf{\defterm{Status update}}}\\
\defterm{Status map} $\chi$ & Live or off, for each region.\\
\defterm{Force flag} & True if the update forces a handoff.\\
\defterm{Target session} $r$ & The session to cut, identified by its index and
session ID $\sigma_r$. Used only if the flag is true.\\
\defterm{Retained height} $h_*$ & The last DATA height that the external chain keeps
from session $r$. Used only if the flag is true. It is a height, not a
block hash.\\
\defterm{Activation height} $\lambda_a$ & The update takes effect once L1 block
$\lambda_a$ is finalized. An update included in L1 block
$\lambda_p$ must have $\lambda_a\ge\lambda_p+N$, where the \defterm{notice period}
$N$ is a number of L1 blocks fixed by the initial configuration. The L1
timestamp $\theta_a$ of block $\lambda_a$ only labels the activation; no
rule waits for a local clock to reach it.\\
\bottomrule
\end{tabularx}
\end{table}

\paragraph{Authority records.}
L1 publishes one authority $\tagg{AUTH}_{r,g}$ for each
session index $r$ and generation $g$. $\tagg{AUTH}_{r,0}$ fixes the
session's voting snapshot $\mathit{cfg}_r$ (members, weights, $F_i$,
$\Delta_i$, and $K$), its quota $\TB_r$, and the leader schedule of
generation~$0$. For $g>0$, only the leader schedule of $\tagg{AUTH}_{r,g}$ is
used; the snapshot and the quota stay those of $\tagg{AUTH}_{r,0}$. The handoff chooses region $i_r$
(\cref{sec:witness}), and the session reads its voters and leaders for
that region. L1 does not see regular boundaries, so records are keyed by
index: a session opened by a Force that reuses index $r$ reads the same
records for its own region.

\paragraph{Status updates.}
Status updates follow these rules.
\begin{itemize}
\item Each update is posted at least a configured notice period before
its activation height, as \cref{tab:l1} specifies. Nodes may prepare during this period. The status map and production authority change at activation.
\item Activation heights increase from one update to the next. Nodes
apply updates in L1 order. The \defterm{effective status map}
$\chi(\lambda)$ at L1 height $\lambda$ is the map of the latest valid
update whose activation height is at most $\lambda$. Such an
update lies in an L1 block at or below $\lambda$. Once height $\lambda$ is
finalized, $\chi(\lambda)$ is therefore fixed, and every node computes
the same map.
\item An update whose force flag is false is \defterm{passive};
it changes only the status map. An update whose force flag is true is a
\defterm{Force}; its status map must mark the target region off.
\item Nodes check each update's authorization and L1 finality.
They ignore an update that is unauthorized, violates the
notice period, or does not increase the activation height. A
finalized record never changes.
\end{itemize}
The committee's rules for submitting updates and Forces are
in \cref{sec:force-stages}.

\paragraph{Force scope and validity.}
A Force $F$ names a target session $\sigma$, with index $r$,
and a retained height $h_*$. Its \defterm{deletion scope} $\Del(F)$ is the
set of external blocks that $F$ removes:
\begin{enumerate}[label=(\roman*)]
\item the DATA blocks of $\sigma$ above height $h_*$; and
\item every block of every session that follows $\sigma$ on its
branch, that is, every session reached through $\sigma$'s regular
boundary (\cref{sec:witness}), including sessions that start after $F$
is posted.
\end{enumerate}
$F$ replaces $\sigma$'s regular boundary by a forced one at $h_*$,
even when $h_*=H_r$. If a block is in $\Del(F)$, so are its descendants.
The rest of this document refers to $\Del(F)$ instead of restating it.
A Force is \defterm{applicable} at a node if its target session is
on the node's current branch.

A node also ignores a Force unless all of the following
hold: its target session is on the node's current branch and has not
been the target of an earlier Force; $h(A_{r,0,0})\le h_*\le H_r$ in the
target session; and every earlier Force activated before this one was
included in L1. These checks need the target's boundary. A node that
does not have it yet keeps the record and checks it when the boundary
arrives (\cref{alg:force}). Under \cref{ass:committee}, every committee
update passes all checks.

\paragraph{Genesis.}
The initial configuration fixes the cyclic region order
$\mathcal L_R$, the genesis block, an initial L1 height
$\lambda_{\rm gen}$, and an initial status map that is effective from
$\lambda_{\rm gen}$. A synthetic genesis boundary opens session $0$:
$\lambda^{\rm open}_0$ is the least finalized L1 height at least
$\lambda_{\rm gen}$ at which some region is live, and $i_0$ is the first
live region in $\mathcal L_R$ under $\chi(\lambda^{\rm open}_0)$. Until
such a height is finalized, this boundary waits, like any other
boundary (\cref{sec:witness}).

\subsection{Stake and faults}
\label{sec:stake}
Member $v$ of region $i$ has voting weight $\omega_{i,v}>0$. Let
$W_i=\sum_{v\in V_i}\omega_{i,v}$. The stake configuration
fixes a bound $F_i<W_i/2$, and Byzantine members of region $i$ hold
total weight at most $F_i$. Every regional certificate needs
matching signatures from distinct members with total weight greater than
$F_i$. So every certificate has a correct signer, and the correct members
alone can form one. This replaces the $f+1$ of $n=2f+1$ in Part~I.

Membership and weights stay fixed for a whole session, including its
recoveries. Slashing and status updates do not change them. A later
session checks old certificates against the old snapshot. The bound
$F_i$ counts every key of the snapshot that is compromised while its
certificates are still accepted. Compromise of old keys after that is
outside this model.

Correct replicas keep their protocol state and committed blocks.
Signatures cannot be forged, and hashes are binding.

\subsection{Assumptions}
\label{sec:assumptions}
\begin{assumption}[Finalized L1 records]
\label{ass:authority}
At every finalized L1 block height, the authority records and the ordered
list of status updates through that height are common knowledge. These
finalized records never change.
\end{assumption}
A replica uses only finalized records. If an authority needed for entry
is not yet finalized, the replica waits for it.

\begin{assumption}[Authenticated synchronous communication and global clock]
\label{ass:network}
\leavevmode\par
\begin{itemize}
\item Messages use authenticated communication channels.
\item Within region $i$, the actual delay of a correct-to-correct
message is $\delta\le\Delta_i$; $\delta$ may differ across messages.
\item All nodes share a global clock.
\end{itemize}
\end{assumption}
The delay bound includes delivery and processing of complete validation
dependencies. Each replica records its own entry and voting events on
the global clock.

\begin{assumption}[Cross-region backbone]
\label{ass:backbone}
A backbone eventually delivers each session's committed blocks, with
their regional receipts, to all other nodes. Its delay has no fixed
bound.
\end{assumption}

\begin{assumption}[Honest multisignature committee]
\label{ass:committee}
The multisignature committee is one honest, coordinated service. It
keeps durable state for validated receipts, their first validation times,
a protected prefix, and its outstanding Force. It follows
\cref{sec:l1,sec:force}. It never authorizes a Force from a stale copy
of its protected prefix.
\end{assumption}

\begin{assumption}[Bounded monitoring and Force finalization]
\label{ass:monitor}
There are finite bounds $T_{\rm mon}$ and $T_{\rm L1}$,
measured on the shared global clock, such that:
\begin{itemize}
\item \emph{Monitoring.} When a correct observer validates and sends a complete regional
finality bundle at time $s$, the committee receives it and validates its historical
regional evidence by $s+T_{\rm mon}$. The bound includes required data, missing history,
and processing. Bundles may be batched and may use resolvable caches.
\item \emph{Force finalization.} A valid Force transaction submitted through the L1 service at
time $s$ becomes finalized and available to correct observers by
$s+T_{\rm L1}$. This bound includes transaction inclusion and finalization.
\end{itemize}
\end{assumption}
The committee's decision window $T_{\rm msig}$ is a configured
parameter, not an environment bound; \cref{sec:protection} gives the rule
that uses it. The monitoring path has a bound; the ordinary backbone need only deliver
eventually. Meeting $T_{\rm mon}$ requires the configured data and
processing capacity, not honesty alone. No monitoring acknowledgement is
required. L1 block timestamps label L1 blocks but do not state when they
became finalized; every deadline in this document uses the global clock.

\subsection{Conditions for progress}
\label{sec:progress}
The deferred liveness analysis will also use these conditions.
\begin{itemize}
\item \textbf{Schedules.} Correct leaders get repeated turns in views
that are actually entered. This includes recovery initializers and
DATA-producing turns across sessions. Turns in views that are never
reached do not count.
\item \textbf{Exclusion.} Leaders exposed for equivocation are eventually
excluded by later authorities, as in Part~I's accountability assumption.
\item \textbf{Committee service.} The committee gives working regions
enough uninterrupted time to produce DATA, complete handoffs, and pass
the global waiting period.
\item \textbf{Handoff service.} Every submitted Force eventually reaches
its finite activation height. Required authorities eventually become
available. A waiting handoff eventually has another live region and
can obtain its anchor data from outside the failed region.
\item \textbf{Requests.} A correct client's request that remains
admissible eventually reaches every region that continues to serve.
Each request fits in one DATA batch. Each node keeps pending requests
in the order it first receives them, with only finitely many requests
ahead of any one request (\cref{sec:requests}).
\end{itemize}

\section{Regional engine: definitions}
\label{sec:regional}

\Cref{tab:notation} collects the notation used below. Anchor indices
are always written in full.
\begin{table}[htbp]\centering\small
\caption{Notation.}\label{tab:notation}
\begin{tabularx}{\linewidth}{@{}p{.29\linewidth}Y@{}}
\toprule Symbol & Meaning\\\midrule
$(r,g,t)$; $j$ & Session, recovery generation, leader term; position
within that term. Indices refer to one branch, on which index $r$ names session $\sigma_r$.\\
$\sigma_r$ & \defterm{Session ID}: the hash of the boundary that opened
session $r$ (\cref{sec:witness}).\\
$i_r$; $V_i$ & Region serving session $r$; members of region $i$.\\
$\omega_{i,v}$; $W_i$; $F_i$ & \defterm{Member weight}; \defterm{total regional weight};
configured upper bound on Byzantine weight, with $F_i<W_i/2$.\\
$\Delta_i$; $K$ & \defterm{Regional message-delay bound}; maximum DATA blocks
in a leader term.\\
$\tagg{AUTH}_{r,g}$; $\mathit{cfg}_r$ & \defterm{Finalized authority} for the
generation; identifier of the session's \defterm{fixed voting snapshot}.\\
$A_{r,g,t}$ & \defterm{Anchor} immediately before term $(r,g,t)$.\\
$A_{r,0,0}$; $A_{r,g,0}$ & \defterm{Inherited session anchor}; \defterm{generation anchor}
chosen at initialization or recovery, also called the generation's
\defterm{opening}.\\
$h(B)$; $k_r(B)$ & \defterm{DATA height} from genesis; \defterm{physical position} within
session $r$, counting DATA and DRAIN.\\
$\TB_r$; $H_r$; $T_r$ & \defterm{Session DATA quota}; \defterm{last allowed DATA height};
DATA block at that height.\\
$m_{r,g,t}$ & Number of fresh positions in the term.\\
$L(r,g,t)$ & Leader of the view.\\
$\QC_{r,g,t}(B)$; ANCHOR QC & \defterm{Proposal certificate} for $B$;
certificate re-certifying $A_{r,g,0}$.\\
$\FC_r(B)$ & \defterm{Regional receipt}: matching FINAL signatures of weight
$>F_{i_r}$, validated by \cref{def:receipt}.\\
$Q_r^0$; $Q_v$; $Q^*$ & Inherited anchor's commitment evidence; replica's \defterm{entry
lock}; justification selected by the recovery leader.\\
$e_v(r,g,t)$; $\tau_v(P)$; $a_v$ & \defterm{Local view-entry time}; time of first
validated vote for proposal $P$; time of receiving a quit proof.\\
$\rho$; $\lambda(B)$ & Block's \defterm{finalized L1 reference}; its L1 height.\\
$\chi(\lambda)$; $\mathcal L_R$ & \defterm{Effective live/off map} at L1 height
$\lambda$; \defterm{cyclic region order}.\\
$h_*$; $\lambda_a$; $N$ & \defterm{Retained DATA height} of a
Force; \defterm{activation L1 height}; \defterm{notice period} in L1 blocks.\\
$\Del(F)$ & \defterm{Deletion scope} of Force $F$ (\cref{sec:l1}).\\
$\lambda_s$; $\lambda_r^{\rm open}$ & \defterm{Successor-selection L1 height};
selection height of the boundary that opened session $r$.\\
$G_u$; $G_{\rm msig}$ & Last block of observer $u$'s
\defterm{global prefix}, initially a trusted checkpoint; the committee's internal
\defterm{protected checkpoint}. Both include session IDs.\\
$s_u(B)$; $\mu(B)$ & First full validation and relay by observer $u$;
first full validation by the committee.\\
$T_{\rm mon}$; $T_{\rm msig}$; $T_{\rm L1}$ & \defterm{Monitoring bound};
\defterm{committee decision window}, a parameter; \defterm{Force submission-to-finalization bound}.\\
$\Twait$ & \defterm{Observer waiting period},
$T_{\rm mon}+T_{\rm msig}+T_{\rm L1}$.\\
$E$; \defterm{Ready} & \defterm{Entry proof}; completion of the
successor's boundary, history, and authority validation (\cref{def:ready}).\\
\bottomrule
\end{tabularx}
\end{table}

\subsection{Blocks, positions, and L1 references}
\label{sec:blocks}
The \defterm{anchor} $A_{r,g,t}$ is the block immediately before the first fresh
position of view $(r,g,t)$. We always write all three indices.
\begin{itemize}
\item $A_{r,0,0}$ is the DATA block inherited from the preceding
session; $A_{0,0,0}$ is genesis. It keeps the regional commitment evidence from the session that
produced it: a direct or covering receipt as defined in
\cref{sec:coverage}. Genesis has synthetic evidence.
\item For $g>0$, NEWVIEW chooses the \defterm{recovery anchor} $A_{r,g,0}$
(\cref{alg:recovery}). It is certified but may be uncommitted.
\item For $t>0$, $A_{r,g,t}$ is the terminal DATA block of term
$(r,g,t-1)$, certified by that term's terminal QC.
\end{itemize}
The anchor of a view never changes. The generation starts at
$A_{r,g,0}$; later terms use their own anchors $A_{r,g,t}$. ANCHOR voting
re-certifies only $A_{r,g,0}$, once per generation.

Each block has two positions.
\begin{itemize}
\item The \defterm{DATA height} $h(B)$ counts the DATA blocks on the chain
from genesis to $B$.
\item The \defterm{physical position} $k_r(B)$ counts DATA and DRAIN blocks
within session $r$, with $k_r(A_{r,0,0})=0$. This local count does not change
the signed fields of $A_{r,0,0}$.
\end{itemize}
A new block of session $r$ has the form
\[
 B=(r,k,h,\mathit{type},x,\Hash(B_{\rm parent}),\rho),
 \qquad \mathit{type}\in\{\tagg{DATA},\tagg{DRAIN}\},
\]
with $k=k_r(B_{\rm parent})+1$. A \defterm{DATA block} carries a batch $x$ of
requests and has $h=h(B_{\rm parent})+1$. A \defterm{DRAIN block} carries no batch
and has $h=h(B_{\rm parent})$.

\paragraph{L1 references.}
The \defterm{reference} $\rho$ names a finalized L1 block. It holds that L1
block's height $\lambda(B)$ and hash, and a registry proof for the L1
control records up to that height. These rules hold:
\begin{itemize}
\item along DATA ancestry, $\lambda$ never decreases;
\item the first DATA block of session $r$ has $\lambda\ge
\lambda^{\rm open}_r$, the selection height of the boundary that opened
session $r$ (\cref{sec:witness}; for session $0$, \cref{sec:l1});
\item a DRAIN block copies the reference of $T_r$.
\end{itemize}
The authenticated timestamp of the L1 block named by $\rho$ can be
read from its header. It is not a commit timestamp for $B$. Recovery
may commit $B$ much later (\cref{sec:late-receipt}). The
regional Vote step applies no receive-time freshness filter to $\rho$.

The reference of $T_r$ gives the starting L1 height for selecting the
regular successor (\cref{sec:witness}). It cannot protect a block from
an activated force, even if the reference predates the force.

Regional validity uses only the session's fixed authority. Later L1
updates do not change it. A regionally valid block can still be left out
of the external chain (\cref{sec:external}).

\subsection{Session quota, term positions, and DRAIN}
\label{sec:quota}
\paragraph{1. Fix the session endpoint.}
Session $r$ may add $\TB_r\ge1$ DATA blocks after $A_{r,0,0}$. Its \defterm{last
DATA height} is
\[
 H_r=h(A_{r,0,0})+\TB_r.
\]
The DATA block at this height is $T_r$. Recovery never changes $H_r$.
The quota counts the DATA blocks after $A_{r,0,0}$ in a proposal's
ancestry, whether committed or not. DRAIN, resends, and ANCHOR votes
use no quota.

\paragraph{2. Produce DATA while space remains.}
If $h(A_{r,g,0})<H_r$, generation $(r,g)$ is a \defterm{DATA generation}.
Each term begins at its anchor $A_{r,g,t}$ and has
\[
 m_{r,g,t}=\min\{K,\;H_r-h(A_{r,g,t})\}
\]
new positions, numbered $j=0,\ldots,m_{r,g,t}-1$. At position $j$, the
DATA block has
\[
 h=h(A_{r,g,t})+j+1,\qquad
 k=k_r(A_{r,g,t})+j+1.
\]
A new position is called a \defterm{fresh position}; its block is a
\defterm{fresh block}. The anchor itself is not fresh. The last position
$j=m_{r,g,t}-1$ is the \defterm{terminal position}, and its QC is the
\defterm{terminal QC}.

If the terminal DATA height is below $H_r$, its QC admits term
$(r,g,t+1)$, with that terminal block as $A_{r,g,t+1}$. If it is $H_r$,
there is no next term. Only the last DATA term can have fewer than $K$
positions. The initial generation is always a DATA generation, since
$\TB_r\ge1$.

\paragraph{3. Use DRAIN when recovery starts at the endpoint.}
If $g>0$ and $h(A_{r,g,0})=H_r$, generation $(r,g)$ is a
\defterm{drain generation}. It has only term $t=0$ and position $j=0$,
with $m_{r,g,0}=1$. Its one fresh block is DRAIN, whose parent is
$A_{r,g,0}$ and whose justification is the generation's ANCHOR QC.
It carries no requests, keeps DATA height $H_r$, and has no next term.

DRAIN supplies a commit timer. ANCHOR voting starts no timer, so
re-certifying $A_{r,g,0}$ alone does not commit it. At height $H_r$, a
DATA child would exceed the quota. The DRAIN child's timer instead
commits its ancestors, including $T_r$.
\begin{itemize}
\item A DATA generation that reaches $H_r$ does not append DRAIN.
It waits for $T_r$'s timer or recovers after a quit.
\item A later drain generation may start from an earlier DRAIN.
Each interrupted recovery may add one DRAIN, so their total number
is not bounded.
\item The choice depends on the recovery anchor's DATA height, not on
whether a replica has already committed $T_r$.
\end{itemize}
DRAIN stays in the regional log. The external chain contains only DATA,
and the next session inherits $T_r$.

\paragraph{Example.}
Let $h(A_{r,0,0})=100$, $\TB_r=10$, and $K=4$, so $H_r=110$.
\begin{center}\small
\begin{tabular}{@{}lll@{}}\toprule
View & Anchor height & Fresh DATA heights\\\midrule
$(r,0,0)$ & 100 & 101--104\\
$(r,0,1)$ & 104 & 105--108\\
$(r,0,2)$ & 108 & 109--110\\\bottomrule
\end{tabular}
\end{center}
Recovery from DATA 107 starts a new generation whose first term produces
108--110. Recovery from certified DATA 110 starts a drain generation.
In both cases, the endpoint remains 110.

\subsection{Certificates and signed messages}
\label{sec:certificates}
\Cref{tab:certificates} lists the certificates of a session. Each needs
matching signatures from distinct members of $V_{i_r}$ with total weight
$>F_{i_r}$.

\begin{table}[htbp]\centering\small
\caption{Certificates of a session.}
\label{tab:certificates}
\begin{tabularx}{\linewidth}{@{}>{\raggedright\arraybackslash}p{.25\linewidth}>{\raggedright\arraybackslash}p{.24\linewidth}Y@{}}
\toprule
Certificate & Signed message & Purpose\\
\midrule
\defterm{Proposal QC} $\QC_{r,g,t}(B)$ & DATA-VOTE or DRAIN-VOTE & Justifies a
child. A terminal QC below $H_r$ also admits the next term.\\
\defterm{ANCHOR QC} & ANCHOR-VOTE & Re-certifies the opening $A_{r,g,0}$ in
generation $g$, before any fresh proposal.\\
\defterm{BLAME proof} & BLAME & Freezes a generation (\cref{sec:quit}).\\
\defterm{Regional receipt} $\FC_r(B)$ & FINAL & Proves that a correct replica
committed DATA block $B$.\\
\bottomrule
\end{tabularx}
\end{table}

A correct replica casts at most one fresh vote per $(r,g,t,j)$ and one
ANCHOR vote per $(r,g)$. DATA and DRAIN votes follow the same rule but
are different signed types. No QC and no ANCHOR vote starts a commit
timer.

\paragraph{Signed context.}
Every signature binds the session ID $\sigma_r$, the
voting configuration, session $r$, region $i_r$, and the message type.
This separates ordinary continuation from a later force fork
(\cref{sec:witness}). Messages of a view also bind the L1
authority and $(r,g,t)$. The other signed fields are:
\begin{itemize}
\item \textbf{PROPOSE and fresh votes:} the opening identity
(below), position $j$, block
type, physical position, DATA height, and block hash;
\item \textbf{NEWVIEW and ANCHOR-VOTE:} the identity of the chosen
opening. NEWVIEW also binds its justification and entry proof;
\item \textbf{STATUS:} the sender's entry lock and entry proof. No
opening is needed;
\item \textbf{BLAME:} the view only. It binds no opening, block hash, or
progress count. All complaints for a view thus sign the same statement,
even before NEWVIEW;
\item \textbf{FINAL:} session-stable fields only (below).
\end{itemize}
Message displays omit these context fields.

\paragraph{Rank and the base justification.}
Within session $r$, a QC formed in term $(r,g,t)$ for block $B$ has \defterm{rank}
$(g,t,k_r(B))$, compared lexicographically. The ANCHOR QC of generation
$g$ has rank $(g,0,k_r(A_{r,g,0}))$, just below the generation's first
fresh QC. A fixed total order on certified content breaks ties. Signer
sets and encodings do not matter. The \defterm{opening
identity} of generation $g$ is
$(g,k_r(A_{r,g,0}),h(A_{r,g,0}),\Hash(A_{r,g,0}))$. In this document,
``opening'' always means the generation's opening block $A_{r,g,0}$, as
in Part~I; a session is named by its session ID.

The anchor $A_{r,0,0}$ also has a \defterm{base justification} $Q_r^0$: the
anchor together with its direct or covering commitment evidence in the
boundary witness (\cref{sec:coverage,sec:witness}). Its justified block
is exactly $A_{r,0,0}$, even if the evidence uses a receipt for a higher
old-session block. That higher block is evidence only. $Q_r^0$ ranks
below every QC of session $r$.
Recovery uses it when no QC of the session is known. A QC of another
session is never a justification in session $r$, so it cannot replace
$A_{r,0,0}$.

\paragraph{Regional receipts.}
After a replica commits a DATA block $B$ of session $r$, by $B$'s own
timer or as an ancestor, it sends
\[
 \langle\tagg{FINAL},\mathit{cfg}_r,r,h(B),\Hash(B)\rangle_v
\]
to its region and to the backbone. Matching FINAL messages of weight
$>F_{i_r}$ form the \defterm{regional receipt} $\FC_r(B)$. The receipt carries $B$ and the ancestry needed to validate it, or
refers to a validated local cache.

\begin{definition}[\defterm{Valid regional receipt}]
\label{def:receipt}
A regional receipt $\FC_r(B)$ is valid when all of the following hold:
\begin{enumerate}[label=(\roman*)]
\item $B$ is a DATA block of session $r$, extends $A_{r,0,0}$, and has
$h(A_{r,0,0})<h(B)\le H_r$;
\item distinct members in session $r$'s fixed voting snapshot have
valid signatures on the same
$\langle\tagg{FINAL},\mathit{cfg}_r,r,h(B),\Hash(B)\rangle$,
with total weight $>F_{i_r}$ and the signed context specified above; and
\item $B$ and its required ancestry pass the block-validation rules of
\cref{sec:admission}.
\end{enumerate}
The verifier checks this evidence; it does not replay the signers'
timers. The weight threshold guarantees at least one correct signer,
who sends FINAL only after committing $B$.
\end{definition}
Here $\mathit{cfg}_r$ identifies session $r$'s \defterm{fixed voting snapshot}.
When we write $\FC(B)$ without a session subscript, we mean the receipt
from the session that produced $B$. The inherited anchor keeps that
receipt; it does not acquire a new one. A verifier outside
the region checks condition~(iii) from evidence, as
\cref{sec:admission} explains.
\begin{itemize}
\item FINAL binds no generation or term. So FINAL messages for one block
combine across recoveries. The block hash keeps the block's signed
context.
\item If a DRAIN commits $T_r$, replicas send FINAL for $T_r$, never for
the DRAIN.
\item Replicas of session $r$ send no FINAL for $A_{r,0,0}$. It keeps its original commitment evidence, which may come from an
earlier session if intermediate sessions added no DATA.
\item A replica may also commit $B$ and its uncommitted
ancestors when it obtains a valid $\FC_r(B)$ of its own session.
Completion does this for $T_r$ (\cref{alg:steady}). Such a commit
sends no FINAL, because the receipt already exists. Only Completion
also cancels timers. Receipts of other sessions serve only as history or anchor
evidence. They never change this session's locks or log.
\end{itemize}

\paragraph{Receipt coverage.}
\label{sec:coverage}
A valid $\FC_r(D)$ also proves commitment of every DATA ancestor $B$
of $D$ in the same session. At least one correct signer committed
$D$, which commits those ancestors. The \defterm{covering evidence} includes the
valid receipt for $D$ and the fully verified ancestry linking $D$ to $B$.
It does not claim that the signatures name $B$ directly.

This coverage never substitutes for a receipt from a different session.
It also does not replace the check that $B$ belongs to the authorized
external branch. A later Force may discard $D$ while keeping $B$; the
regional commitment evidence for $B$ still holds. External observers,
retained-anchor witnesses, and base justifications may use direct or
covering evidence. No extra FINAL messages or certificate round are
needed. We use ``\defterm{commitment evidence} for $B$'' for either form.

\subsection{Admission, validation, and engine states}
\label{sec:admission}
\paragraph{\defterm{Admission.}}
Replica $v$ admits views as follows.
\begin{itemize}
\item $(r,0,0)$ is admitted at $v$'s Ready event for the boundary that
opens session $r$ (\cref{def:ready}). The witness is the \defterm{entry proof}.
The NEWVIEW must choose exactly $A_{r,0,0}$.
\item $(r,g,0)$ with $g>0$ is admitted once $v$ holds
$\tagg{AUTH}_{r,g}$ and a quit proof for generation $g-1$. That proof is
the \defterm{entry proof}. The NEWVIEW is justified by a QC of an earlier
generation of session $r$, or by $Q_r^0$.
\item $(r,g,t)$ with $t>0$ is admitted once $v$ also holds the terminal
QCs of terms $0,\ldots,t-1$, each for DATA below $H_r$.
\end{itemize}
Term zero is admitted before its NEWVIEW arrives. Its progress checks
therefore run even if the initial leader is silent. Fresh proposals need
the generation's ANCHOR QC. This holds in every generation, including
generation $0$ of session $0$; unlike Part~I, genesis has no shortcut. A
message is \defterm{admissible} if its view is admitted and its signatures
are valid.

A verifier that is not a replica of the session has no
admission events. When it validates history (\cref{def:receipt},
\cref{sec:external}), a view counts as admitted if the evidence contains
its \defterm{admission witness}: the boundary witness for $(r,0,0)$, the chain of
entry proofs for $(r,g,0)$, and the earlier terminal QCs for $t>0$.

\paragraph{\defterm{Complete validation.}}
Validation checks signatures, authority, admission, block type, quota,
both positions, the L1-reference rules, the parent rule
(\cref{tab:parent}), parent QCs, hash links, and application validity
throughout the required ancestry. Every relevant ancestor includes its
signed proposal and block data; a chain of certificate hashes alone is
not enough. Dependencies come from the message or a validated local
cache. Correct senders send \defterm{complete validation bundles}. A reference is
allowed only if the receiver can resolve it within the delay bound of
\cref{ass:network}. Bundles list each dependency once.

\paragraph{Evidence takes priority.}
Before validating a message or discarding a duplicate, inspect all signed
material in it for a quit proof (\cref{sec:quit}). A proof needs its own
valid signatures and admission witnesses. A valid proof counts even if
the message that carries it is invalid or incomplete. Compare incoming
material with all retained material, not only with objects inside the
message.

\paragraph{One proposal-processing rule.}
Every admitted fresh proposal of the current generation that $v$ has not
voted for invokes the Vote step. This includes proposals in attached
histories, proposals of earlier terms, and proposals whose blocks are
already committed as ancestors. A proposal's identity is
$(r,g,t,j,\Hash(B))$. The signer set or encoding of its parent QC does
not change it. Deduplicate by this identity after evidence inspection.
Old-generation proposals are retained without new votes or timers.
Future-generation messages wait for entry.

\paragraph{Engine states.}
A replica runs one engine for each session ID of its region.
Engines of different sessions remain separate, even if they use the
same session index. The engine of
session $r$ is in one of three states.
\begin{itemize}
\item \textbf{\defterm{Active}:} the current generation is admitted and unfrozen.
The engine proposes, votes, and changes terms.
\item \textbf{\defterm{Frozen}:} the current generation has a quit proof. The
engine stops proposals, votes, and term changes. It still processes
evidence, old votes, QCs, receipts, and recovery.
\item \textbf{\defterm{Complete}:} the replica holds $\FC_r(T_r)$ satisfying
\cref{def:receipt}.
Production in session $r$ is over, even if the external handoff is still
pending. The engine only serves data.
\end{itemize}
Committing $T_r$ locally does not complete the engine. Other members may
still need this replica's recovery votes before a receipt forms. Each
message goes to the engine identified by its session ID, so
evidence from one session or branch never freezes another. These engine
states do not describe external finality. The external chain may move
to another branch while an old engine still runs. L1 status
updates and Forces never change an engine's state, timers, or log; they
act only on the external branch (\cref{alg:rotation}).

\begin{table}[htbp]\centering\small
\caption{Parent selection for the next fresh position $j$ of term $(r,g,t)$.}
\label{tab:parent}
\begin{tabularx}{\linewidth}{@{}l Y Y@{}}
\toprule
Condition & $B_{\mathrm{parent}}$ & $Q_{\mathrm{parent}}$\\
\midrule
$j=0,t=0$ (DATA or DRAIN) & Opening $A_{r,g,0}$ & Its ANCHOR QC in $(r,g,0)$\\
$j=0,t>0$ & Anchor $A_{r,g,t}$, the terminal block of $(r,g,t-1)$ & Its proposal QC in $(r,g,t-1)$\\
$1\le j<m_{r,g,t}$ & Position $j-1$ of $(r,g,t)$ & Its proposal QC in $(r,g,t)$\\
\bottomrule
\end{tabularx}
\end{table}

\section{Regional protocols}
\label{sec:protocols}
Protocols 1 and 2 run inside the engine of one session $r$. They are
Protocols 1 and 2 of Part~I~\cite{part1} with the changes listed below.
\RR{Red text} marks every difference from Part~I, in the boxes, in the
message-processing order, and in the quit proofs. Explanations stay
outside the boxes.

\paragraph{Notation.}
All anchors and view times use their full session indices. The displayed
regional messages omit $r$, but their signatures bind the session ID
$\sigma_r$ and $r$. In these two protocols, $\Delta$ and $F$ mean
$\Delta_{i_r}$ and $F_{i_r}$.

\paragraph{Changes from Part~I.}
\begin{enumerate}[label=\arabic*.]
\item \emph{Stake.} Certificates need weight $>F$, not $f+1$ signatures.
\item \emph{Blocks.} A block also carries its session, physical
position, type, and L1 reference $\rho$ (\cref{sec:blocks}). Validity
checks all of them.
\item \emph{Quota.} A term has $m_{r,g,t}\le K$ positions. No term follows
$H_r$.
\item \emph{DRAIN.} A drain generation proposes one DRAIN. Its vote
counts as progress.
\item \emph{Receipts.} Replicas send FINAL after committing DATA. A
receipt for $T_r$ completes the engine (new step~5 of \cref{alg:steady}).
\item \emph{Timers.} Completion also cancels timers. A force never does.
\item \emph{Recovery entry.} After the $\Delta$ wait, a replica enters
when $\tagg{AUTH}_{r,g+1}$ is finalized. $Q_r^0$ can serve as the entry
lock and as the justification.
\item \emph{Genesis.} Session $0$ also certifies its anchor with an
ANCHOR QC before its first fresh proposal.
\item \emph{Sessions.} Each message goes to the engine of its session.
Session $0$ opens by bootstrap (\cref{alg:bootstrap}).
\end{enumerate}

\paragraph{Message-processing order.}\leavevmode
\RR{Replica $v$ first passes each message to the engine of the session
ID that the message binds. L1 updates go to \cref{alg:force}. Boundary witnesses and receipts of
other sessions go to \cref{alg:rotation,alg:external}.} The engine handles each received
message in this order:
\begin{enumerate}[label=(\arabic*)]
\item check all signed material for a quit proof;
\item check completeness and validity;
\item process attached DATA \RR{and DRAIN} proposals under Vote, parents
first;
\item process attached terminal QCs under step~4 of \cref{alg:steady},
in increasing term order;
\item process the message itself.
\end{enumerate}
When a message and a timer fall at the same time, $v$ handles the
message first.

\subsection{Steady state}
\label{sec:steady}
\Cref{alg:steady} gives the steady state of a session: proposals, votes,
timers, term rotation, and completion.

\begin{protocol}{alg:steady}{Propose, vote, commit, normal rotation\RR{, and completion}}
Replica $v$ runs these steps \RR{in the engine of session $r$}. Its
current term is $(r,g,t)$.
\begin{enumerate}[label=\textbf{\arabic*)},itemsep=3pt]
\item \textbf{Propose.} Run this step only if $v=L(r,g,t)$, $g$ is
unfrozen\RR{, and session $r$ is not complete}.
\RR{In a DATA generation,} for each position
$j=0,\ldots,\RR{m_{r,g,t}}-1$, in order:
 \begin{itemize}
 \item Wait until $v$ holds $B_{\mathrm{parent}}$ and
       $Q_{\mathrm{parent}}$ for $j$ (\cref{tab:parent}).
 \item Build \RR{$B=(r,k,h,\tagg{DATA},x,\Hash(B_{\mathrm{parent}}),\rho)$,
       with $k$ and $h$ from \cref{sec:quota} and an L1 reference $\rho$}.
       The batch $x$ follows \cref{sec:requests}.
 \item Broadcast
       $P=\langle\tagg{PROPOSE},g,t,j,B,Q_{\mathrm{parent}}\rangle_v$
       with its validation bundle \RR{to regional replica and the backbone. }
 \item Process $P$ under Vote.
 \end{itemize}
 \RR{In a drain generation, propose once, at $(t,j)=(0,0)$. The block is
 a DRAIN with parent $A_{r,g,0}$, no batch, DATA height $H_r$, and the L1
 reference of $T_r$. Broadcast and process it as above.}
 Do not wait for the parent to commit. Never sign two different
 blocks for one position. A resend repeats the same $P$.

\item \textbf{Vote.} Run this step on every DATA \RR{or DRAIN} proposal
$P$ of term $(r,g_P,t_P)$ at position $j_P$. Here $t_P$ may be earlier
than $t$. The step also applies to proposals inside attached histories.
 \begin{enumerate}[label=(\alph*)]
 \item If $P$ or its attached material contains a quit proof for $g$,
       run Quit and stop.
 \item If $g_P<g$, retain $P$ and stop. If $g_P>g$, buffer $P$
       until $v$ enters generation $g_P$.
 \item If $g$ is frozen \RR{or session $r$ is complete}, retain $P$ and
       stop.
 \item If a dependency is missing, retain $P$, request the dependency,
       and stop. Rerun this step when it arrives.
 \item If $P$ is invalid, discard it and stop. \RR{Validity includes the
       block type, the quota, both positions, and the L1 reference.}
 \item If $v$ has already voted for $P$, stop.
       A duplicate never restarts a timer.
 \item Otherwise, let $\tau_v(P)$ be the current time. Forward $P$ and
       its bundle to all \RR{regional replicas and the backbone. }
       Broadcast
       $\langle\tagg{DATA\mbox{-}VOTE},g,t_P,j_P,\Hash(B)\rangle_v$
       \RR{(DRAIN-VOTE if $B$ is a DRAIN)}.
       Start the commit timer of $P$. It expires at $\tau_v(P)+2\Delta$.
 \end{enumerate}

\item \textbf{Commit.} When the timer of $P$ expires, commit $B$ and
its uncommitted ancestors. \RR{For each DATA block $D$ of session $r$
that this commits, send
$\langle\tagg{FINAL},\mathit{cfg}_r,r,h(D),\Hash(D)\rangle_v$ to the region
and the backbone. A committed DRAIN yields FINAL for $T_r$, never for
itself. Relay each receipt of session $r$ that $v$ obtains to the region
and the backbone.} Only Quit \RR{and Completion} cancel a timer. Entering a
new term does not cancel it. A QC never starts a timer. \RR{A force
instruction never cancels a timer.}

\item \textbf{Terminal QC and rotation.} Run this step the first time
$v$ obtains a valid terminal QC for a term $(r,g,t')$.
``Obtains'' means that $v$ assembles the QC from votes or finds it in
any message, including as another proposal's $Q_{\mathrm{parent}}$.
 \begin{enumerate}[label=(\alph*)]
 \item Process each unseen proposal in the QC's history under Vote,
       parents first.
 \item If $g$ is now frozen, stop. \RR{If the QC certifies DATA at
       $H_r$, or a DRAIN, stop: no next term exists.}
 \item Send the QC and its validation bundle to $L(r,g,t'+1)$.
 \item If $t\le t'$, enter $(r,g,t'+1)$ and set $e_v(r,g,t'+1)$ to the
       current time. Cancel the progress checks of all earlier
       terms.
 \end{enumerate}
 A QC never moves $v$ to an earlier term. It never changes a recorded
 entry time.

\item \RR{\textbf{Completion.} Run this step the first time $v$ obtains a
regional receipt $\FC_r(T_r)$ satisfying \cref{def:receipt}.}
 \begin{itemize}\color{rrRed}
 \item Commit $T_r$ and its uncommitted ancestors.
 \item Mark session $r$ complete. Stop its proposals, votes, progress
       checks, term changes, and recovery. Cancel its unexpired timers.
 \item Relay $\FC_r(T_r)$ to the backbone.
 \item Keep serving data. \Cref{alg:rotation} runs the handoff.
 \end{itemize}
 \RR{A force instruction never runs this step.}
\end{enumerate}
\end{protocol}

The \defterm{acceptance time} $\tau_v(P)$ is when complete validation and the first
vote occur, not an earlier arrival of an incomplete packet. Each replica
records its own entry time for each term. In session $r$, $e_v(r,0,0)$ is
set at Ready, in step~1 of \cref{alg:bootstrap}. For $g>0$, $e_v(r,g,0)$ is
set in step~3 of \cref{alg:recovery}. For $t>0$, $e_v(r,g,t)$ is set in
step~4 of \cref{alg:steady}. A term's progress checks therefore start
when the term is admitted locally, not when its leader's first proposal
arrives. A silent new leader cannot delay them.

A proposal or QC at height $H_r$ does not end the session. Without a
quit, the timer of $T_r$ commits it; no later block is needed. Only a
receipt for $T_r$ completes the engine.

\subsection{Quit proofs and recovery}
\label{sec:quit}
A quit proof freezes a generation. \Cref{alg:recovery} then moves the
session to the next generation.

\begin{protocol}{alg:recovery}{Blame, quit, and recovery}
Replica $v$ runs these steps \RR{in the engine of session $r$}. Its
current term is $(r,g,t)$.
\begin{enumerate}[label=\textbf{\arabic*)},itemsep=3pt]
\item \textbf{Blame.} \RR{Let $q=m_{r,g,t}$ in a DATA generation and
$q=1$ in a drain generation.} For each $p=1,\ldots,\RR{q}$, run a
progress check at time
\[
 e_v(r,g,t)+(2p+4)\Delta.
\]
\RR{Schedule the check for $p=1$ at entry, even if the opening is not yet
known. Schedule the others once $q$ is known. A check whose time has
already passed runs at once.}
At that time, broadcast $\langle\tagg{BLAME},g,t\rangle_v$ if all of
these hold:
 \begin{itemize}
 \item $v$'s current term is still $(r,g,t)$, $g$ is unfrozen\RR{, and
       session $r$ is not complete};
 \item $v$ has cast DATA \RR{or DRAIN} votes for fewer than $p$ positions
       of term $(r,g,t)$ (ANCHOR votes do not count);
 \item $v$ has not yet sent BLAME for $(r,g,t)$.
 \end{itemize}

\item \textbf{Quit.} Run this step the first time $v$ holds a quit proof
for its current generation $g$.
 \begin{itemize}
 \item Freeze $g$. Stop all progress checks, proposals, votes, and
       term changes in $g$. Cancel every unexpired commit timer of $g$.
 \item Set $a_v$ to the current time. Broadcast the proof.
 \item Keep all commitments. Keep collecting votes and QCs of $g$.
       They can raise the entry lock chosen in step~3. \RR{Keep
       processing regional receipts (step~5 of \cref{alg:steady}).}
 \item If the proof is of type (b) or (c), submit it to L1. Do
       not wait for L1 to process it.
 \end{itemize}

\item \textbf{Status and entry.} \RR{Skip this step if session $r$ is
complete.} \RR{Wait until at least $a_v+\Delta$ and until
$\tagg{AUTH}_{r,g+1}$ is finalized. Set $e_v(r,g+1,0)$ to the time both
conditions hold, and read that authority from L1.} Then:
 \begin{itemize}
 \item Set the entry lock $Q_v$ to the highest-ranked QC
       (DATA\RR{, DRAIN,} or ANCHOR) that $v$ knows from generations
       $0,\ldots,g$\RR{, or to $Q_r^0$ if it knows none}.
       $Q_v$ never changes after this. QCs that arrive later do not
       update it.
 \item Enter $(r,g+1,0)$ as unfrozen and start its progress checks.
       \RR{Keep $A_{r,0,0}$, $H_r$, and the voting snapshot of session $r$.}
 \item Send $\langle\tagg{STATUS},g+1,0,Q_v,E\rangle_v$ and its
       validation bundle to $L(r,g+1,0)$. Here $E$ is the quit proof for
       $g$, which is also the entry proof of $g+1$.
 \item Process any buffered proofs for $g+1$.
 \end{itemize}

\item \textbf{New-view.} Run this step only at $L(r,g+1,0)$. Wait $2\Delta$
after its own entry. If $g+1$ is still current and unfrozen:
 \begin{itemize}
 \item Let $Q^*$ be the highest-ranked QC from generations $0,\ldots,g$
       that the leader knows, including the QCs in STATUS
       messages\RR{; if it knows none, let $Q^*=Q_r^0$}.
 \item Let $A_{r,g+1,0}$ be the block that $Q^*$ certifies. \RR{If
       $h(A_{r,g+1,0})<H_r$, generation $g+1$ is a DATA generation.
       Otherwise it is a drain generation.}
 \item Broadcast
       $\langle\tagg{NEWVIEW},g+1,0,A_{r,g+1,0},Q^*,E\rangle_{L(r,g+1,0)}$
       with its validation bundle. Choose only one opening block for
       $g+1$.
 \end{itemize}

\item \textbf{Anchor vote.} On receiving a NEWVIEW for $g+1$:
 \begin{enumerate}[label=(\alph*)]
 \item If it contains a quit proof for $g+1$, run Quit and stop.
 \item Validate its leader, entry proof $E$, opening block $A_{r,g+1,0}$,
       $Q^*$, and dependencies. If data are missing, request them and
       wait without voting. If the message is invalid, discard it.
 \item Vote only if all of these hold: $g+1$ is current and unfrozen; $v$
       has cast no ANCHOR vote in $g+1$; and $\rank(Q^*)\ge\rank(Q_v)$.
 \item To vote, forward the NEWVIEW to all replicas and broadcast
       $\langle\tagg{ANCHOR\mbox{-}VOTE},g+1,0,\Hash(A_{r,g+1,0})\rangle_v$.
       Start no timer. Use no fresh position.
 \end{enumerate}
 Matching ANCHOR votes \RR{of total weight $>F$} form the ANCHOR QC.
 When $L(r,g+1,0)$ obtains it, the leader starts Propose in
 \cref{alg:steady} at $j=0$\RR{; in a drain generation, this is the
 DRAIN}. Every replica, including one that cast no ANCHOR vote,
 processes DATA \RR{and DRAIN} proposals backed by a valid ANCHOR QC
 under Vote.
\end{enumerate}
\end{protocol}

\paragraph{\defterm{Complete quit proofs.}}
For an admitted generation $g$ \RR{of session $r$}, a proof consists of
either:
\begin{enumerate}[label=(\alph*)]
\item distinct signatures \RR{of total weight $>F$} on the same
$\langle\tagg{BLAME},g,t\rangle$ for an admitted term, with that term's
admission witness;
\item two different admissible DATA \RR{or DRAIN} blocks signed by
$L(r,g,t)$ at the same $(r,g,t,j)$, with their admission \RR{and quota}
witnesses; or
\item two valid NEWVIEWs signed by $L(r,g,0)$ selecting different opening
blocks for generation $g$.
\end{enumerate}
\RR{A DATA block beyond the quota, or a DRAIN outside a drain generation,
is invalid. It is rejected and never counts for (b).}
Different openings count even when their blocks are comparable, because
they fix different starting heights $h(A_{r,g,0})$. The height formula then
maps the same position to different heights.

Different QCs for the same opening block are not equivocation. Evidence
from any admitted term can freeze its generation, including a term not
yet locally reached. Old-generation proofs are archival and may be
submitted to L1. \RR{Evidence from another session never freezes this
engine.} Future-generation proofs are buffered until entry. Neither
resets an unrelated recovery.

The entry-lock test governs ANCHOR voting only. A replica that has
committed $T_r$ locally still takes part in recovery until it holds
$\FC_r(T_r)$, because other members may need its votes.

\section{Regional and global finality}
\label{sec:external}
This section uses the branch, session ID, and boundary
witness of \cref{sec:witness}, and the deletion scope $\Del(F)$ of
\cref{sec:l1}.

\subsection{Regional finality}
\label{sec:regional-finality}
Each observer $u$ keeps a \defterm{global prefix}: a sequence
of DATA blocks that $u$ never rolls back. Its last block, with its
session ID, is $u$'s \defterm{global checkpoint} $G_u$. Initially, $G_u$ is
genesis or a checkpoint that $u$ explicitly trusts.
The checkpoint includes, or authenticates, its session ID, rotation
position, voting configuration, and processed L1 history. The observer
uses the current finalized L1 records to extend that history.

\begin{definition}[\defterm{Regional finality}]
\label{def:regional-finality}
A DATA block $B$ of session $\sigma_r$ has regional finality in an observer's current snapshot
when all of the following hold:
\begin{enumerate}[label=(\roman*)]
\item Starting from $G_u$, the regular boundaries and
activated Forces give a current branch on which $\sigma_r$ is the
session that covers DATA height $h(B)$, and $i_r$ is the region of
$\sigma_r$.
% \item The proposal's view $(r,g,t)$, leader, anchor, and signed context are valid under that session's historical authority.
\item $B$ has valid regional commitment evidence: its direct receipt
under \cref{def:receipt}, or a covering receipt under \cref{sec:coverage}.
\item The DATA history from $G_u$ through $B$ is a valid, continuous
extension. Every DATA block has direct or covering commitment evidence
from its own session. No
activated Force removes any part of this history.
\end{enumerate}
\end{definition}
The assigned region is derived before its receipt is checked. A region's
valid receipt does not authorize that region to produce another region's
blocks. The current live/off map also does not replace the historical
authority of a retained block.

Regional finality is relative to the current external branch. A later
Force may remove that output. Its original regional commitment and
receipt remain valid evidence about the old engine. A finalized Force
that has not activated does not yet change this branch; it affects
global finality as specified below. The DATA blocks with
regional finality on $u$'s current branch form $u$'s \defterm{regional
output}.

\paragraph{Remark (Commitment by one replica and by all replicas).}
Here ``all replicas'' means all correct voters of the session that
produced $B$. Consider that session's regional engine, with recovery
authorities set in advance.
\begin{enumerate}[label=(\arabic*),itemsep=4pt]
\item \textbf{What a receipt proves.} A valid direct or covering
receipt proves that at least one correct replica has committed $B$.
Other correct replicas may still have running timers, or may have
canceled their timers on Quit. Regional finality does not assert that
they have all committed already.

\item \textbf{When all replicas commit.} Regional safety preserves
$B$ through recovery: the selected anchor is $B$ or a descendant of
$B$. ANCHOR voting re-certifies that block but starts no commit timer.
Suppose a correct replica votes for and forwards a fresh DATA or DRAIN
proposal for a descendant $D$ at time $s$. If no correct replica freezes
that generation through $s+3\Delta_{i_r}$, all correct replicas commit
$B$ by that time. Replicas that have not completed the session vote by
$s+\Delta_{i_r}$; their timers commit $D$ and its ancestors at most
$2\Delta_{i_r}$ later. Completion may cancel a timer, but first commits
$T_r$ and its ancestors, including $B$. In particular, a DRAIN need not
itself commit at a replica that completes through the receipt for $T_r$.

\item \textbf{Why one honest leader is not enough.} An honest leader
may finish its term before its timers expire. A Byzantine successor
can then equivocate and freeze the generation. What is needed is the
uninterrupted commit window in (2). Under regional safety and the
progress conditions of \cref{sec:progress}, including permanent
equivocation exclusion and fair leader schedules, the engine eventually
obtains such a window while it continues to run, or completes through
the terminal receipt. Earlier recoveries do not revoke $B$. These
conditions do not give a fixed bound on when the window begins.

\item \textbf{When a replica can know.} A correct regional replica
that directly commits $D$ by its timer at time $c$ can conclude that
all correct regional replicas have committed $B$ if it remains Active
in the same generation through $c+2\Delta_{i_r}$, after processing
messages due at that time. Other replicas' timers are due by
$c+\Delta_{i_r}$. A Quit that canceled one of those timers would reach
this replica by $c+2\Delta_{i_r}$ and freeze its generation. A replica
whose timer was canceled by Completion has already committed $B$.
Normal term rotation does not affect this test.

\item \textbf{What an external observer can know.} A receipt alone
does not establish that uninterrupted window. The backbone has no
fixed delay bound, so receiving no quit proof is not evidence that no
replica quit. Thus receipt arrival starts no fixed wait that certifies
commitment by all regional replicas. The global-finality wait defined
next protects the external prefix against a Force; it does not certify
that every regional replica has completed its local commitment.
\end{enumerate}

\subsection{Global finality}
\label{sec:global-finality}
A \defterm{complete regional-finality bundle} contains the blocks, receipts,
application history, boundary witnesses, and finalized L1 proofs needed
to verify the path from a trusted checkpoint. Previously validated data
may be cached. A valid descendant receipt can cover DATA ancestors in
its own session, with their fully verified ancestry. Hashes
alone do not supply missing block data or another session's evidence.

Let $s_u(B)$ be the first time observer $u$ completely validates regional
finality for $B$ and sends its complete bundle to the committee.
Process ancestors before children and do this for every newly
validated DATA block. Set each time once. Different signer sets,
certificate encodings, duplicate bundles, and later recovery do not reset
it. If a node loses its timing record, it waits again from a new relay.
It cannot reconstruct an earlier start from the block's L1 reference.

The \defterm{observer waiting period} is
\[
 \Twait=T_{\rm mon}+T_{\rm msig}+T_{\rm L1}.
\]
The three terms pay for complete monitoring, the committee's decision,
and L1 finalization. The bounds are those of \cref{ass:monitor}, and
$T_{\rm msig}$ is the window of \cref{sec:protection}.

\begin{definition}[\defterm{Global finality at an observer}]
\label{def:global-finality}
At global-clock time $z$, observer $u$ may append DATA block $B$ to its
global prefix if:
\begin{enumerate}[label=(\roman*)]
\item $B$ has regional finality under $u$'s finalized L1
snapshot at time $z$ (\cref{alg:external}, step~1);
\item all preceding DATA blocks on $u$'s current branch are already
in $u$'s global prefix;
\item $z\ge s_u(B)+\Twait$; and
\item $B\notin\Del(F)$ for every L1-finalized Force $F$
whose target session is on $u$'s current branch, whether or not $F$ has
activated.
\end{enumerate}
At a deadline shared with L1 finalization, process the finalized Force
first, then test these conditions.
\end{definition}
Only a continuous DATA prefix becomes global. DRAIN has no external
finality of its own. Regular handoff and regional production do not
wait for this promotion.

This is an observer's verified, persistent state. Observers may confirm
the same block at different times. There is no new global certificate,
per-block committee acknowledgement,
or L1 transaction for promotion. A new observer must
start from a checkpoint it explicitly trusts, or verify the required
history and complete its own waiting periods. A block and its regional
receipt alone do not prove a third party's global-finality claim.

\subsection{Why the block's old timestamp is not the waiting start}
\label{sec:late-receipt}
The L1 reference in a DATA block is fixed when the block is
proposed. The block may obtain a QC, survive several interrupted
recoveries, and commit much later, for example by a DRAIN that keeps the
reference unchanged. A wait measured from that reference's timestamp
could therefore end before any commit evidence exists. The external wait
starts instead at $s_u(B)$, after complete receipt validation and relay;
an old proposal or QC starts no such wait. This affects only observers:
regional proposal timers, locks, and the FINAL format are unchanged.

\section{Cross-region rotation}
\label{sec:xprotocols}

\subsection{Anchors, successor selection, and boundary identity}
\label{sec:witness}
Let $\tagg{\defterm{NextLive}}(i,\chi)$ be the first other live region after $i$
in the cyclic list $\mathcal L_R$. Search at most one cycle. If none
exists, the result is $\bot$ and the handoff waits.

A regular handoff inherits $A_{r+1,0,0}=T_r$. A forced handoff inherits
the DATA block at $h_*$ on the target session's branch. If
$h_*=h(A_{r,0,0})$, it reuses $A_{r,0,0}$ and its original direct or covering commitment
evidence.
\paragraph{\defterm{Successor-selection height.}}
The height $\lambda_s$ identifies the finalized L1 status map used to
choose the next region. It is an L1 block height, not a clock time or
the latest L1 height when a node receives the handoff evidence.
First define the \defterm{starting selection height}
\[
 \lambda_0=\begin{cases}
 \lambda(T_r),&\text{regular handoff},\\
 \max\{\lambda_a,\lambda(A_{r+1,0,0})\},&\text{forced handoff}.
 \end{cases}
\]
For a node evaluating this boundary, let $\ell$ be the latest finalized
L1 height available to it. The node reads the complete status-update
history through $\ell$ and forms the \defterm{set of eligible heights}
\[
 \mathcal S_r(\ell)=
 \bigl\{\lambda\in\mathbb Z_{\ge0}:
 \lambda_0\le\lambda\le\ell,\quad
 \tagg{NextLive}(i_r,\chi(\lambda))\ne\bot\bigr\}.
\]
If $\mathcal S_r(\ell)$ is empty, retain the anchor and wait for further
L1 finalization. Otherwise set
\[
 \lambda_s=\min\mathcal S_r(\ell).
\]
Thus $\lambda_s$ is the first L1 height at or above $\lambda_0$ whose
effective status map has a live region other than $i_r$. Use that map
to select the successor and set
\[
 i_{r+1}=\tagg{NextLive}(i_r,\chi(\lambda_s)),\qquad
 \lambda^{\rm open}_{r+1}=\lambda_s.
\]
By \cref{sec:l1}, finalized heights and their effective status maps
never change. A larger $\ell$ only adds candidate heights above the
previous $\ell$, so it cannot change a minimum already found. Nodes
that can select therefore obtain the same $\lambda_s$ and successor;
a node with insufficient finalized history waits. Later passive status
updates do not reselect this boundary.
An applicable activated Force may replace a regular boundary.
A forced boundary is never replaced, because a session is
the target of at most one Force. It is discarded only together with its
branch, by a Force on an earlier session.

\paragraph{Example: regular selection after a wait.}
Let the cyclic region order be $A\to B\to C\to A$, let $i_r=A$,
and suppose $\lambda(T_r)=100$, so $\lambda_0=100$. Regions $B$ and
$C$ are off at every L1 height from 100 through 104. A passive update
makes $B$ live at activation height 105. A node whose latest finalized
height is 104 has no eligible height and waits. Once height 105 is
finalized, it selects $\lambda_s=105$ and $i_{r+1}=B$.
Suppose another passive update makes $B$ off and $C$ live at height
120. A node that first receives the terminal receipt at finalized
height 120 still selects height 105 and region $B$: it uses the earliest
eligible historical map, not the map at receipt time.

A \defterm{boundary witness} contains its REGULAR or FORCE \defterm{mode}, the predecessor
session ID, the retained anchor and its commitment evidence, the successor,
$\lambda_s$, and the needed L1 proofs. A \defterm{REGULAR witness}
carries the predecessor's \defterm{terminal receipt}, a valid
$\FC_r(T_r)$. A \defterm{FORCE witness} also identifies
the activated instruction. The successor's finalized authority must be
available before entry. Genesis uses a synthetic witness and receipt
(\cref{sec:l1}).

\paragraph{Session IDs and branches.}
The boundary's \defterm{handoff type} is
$\text{\defterm{mode}}\in\{\tagg{REGULAR},\tagg{FORCE}\}$:
\tagg{REGULAR} opens the successor after the predecessor completes its
quota, using $T_r$ as the inherited anchor; \tagg{FORCE} opens it through
an activated L1 Force, using that Force's retained anchor.
The \defterm{canonical identity} of a boundary determines the session
ID of the session it opens:
\[
 \sigma_{r+1}=\Hash\bigl(\sigma_r,\ \text{\defterm{mode}},\ \Hash(A_{r+1,0,0}),\
 i_{r+1},\ \lambda_s,\ \text{\defterm{Force identity}, if any}\bigr).
\]
The handoff type distinguishes a forced successor from a regular one,
even if both inherit the same anchor and select the same region.
Signer sets and encodings do not change the session ID. Regional signatures already
bind $\sigma_r$ (\cref{sec:certificates}). A \defterm{branch}
is a chain of sessions $\sigma_0,\sigma_1,\ldots$ in which each witness
names the previous session. A node's \defterm{current branch} starts at its
checkpoint and follows the regular boundaries, except where an activated
Force has replaced one (\cref{alg:rotation}). An \defterm{external block} is
identified by the pair $(\sigma_r,\Hash(B))$. Protected and global
prefixes retain this provenance; a matching index, height, or
hash in another session is not a replacement for the same block.
Receipts, records, and timers are keyed by this pair.

\subsection{Regular and forced handoff}
\begin{protocol}{alg:rotation}{Rotate Regions Fork Choice Rule.}
Each node runs this handler on its current branch.
\begin{enumerate}[label=\textbf{\arabic*)},itemsep=3pt]
\item \textbf{Regular completion.} On a valid terminal receipt
$\FC_r(T_r)$:
\begin{itemize}
\item Verify the authorized regional source and exactly $\TB_r$ DATA
blocks after $A_{r,0,0}$.
\item Apply current activated Forces. If they remove this regular
boundary, retain its receipt only as old regional evidence.
\item Otherwise select the successor under \cref{sec:witness} and
form the REGULAR witness. Send the witness and its full
bundle to the members of $V_{i_{r+1}}$ over the backbone.
\end{itemize}
Do not wait for $T_r$ to become global. A submitted or finalized but
inactive Force does not stop this regular path.

\item \textbf{Activated Force.} On notification from \cref{alg:force}:
\begin{itemize}
\item Remove $\Del(F)$ from the current branch: the target
session's DATA above $h_*$ and every session after its replaced regular
boundary. This includes regular successors that started during the notice period.
Archive their evidence.
\item Keep the global prefix $G_u$ unchanged, with the same
block provenance. A valid Force cannot remove it.
\item Fetch the retained anchor, its original commitment evidence, and the missing
application history. If they are not locally available, wait.
\item Select the successor and form the FORCE witness. Send
it with its full bundle to the members of $V_{i_{r+1}}$, which start the
new session when Ready.
\end{itemize}
The old regional engines keep their own logs and protocol state. This
handler does not cancel their commit timers or create a quit proof.

\item \textbf{Track the branch.} Use the session ID to route
messages and store receipts. A new Force session may reuse a numeric
session index, but it never reuses the old session's locks, votes,
entry clock, or observer waiting times.
\end{enumerate}
\end{protocol}

\subsection{Ready and the common bootstrap}
\label{sec:ready}
\begin{definition}[\defterm{Ready}]
\label{def:ready}
A member $v$ of the chosen successor is \defterm{Ready} when it has completely
validated a boundary witness on its current branch, its
anchor and missing application history, and the matching finalized
authority.
A regular witness proves the terminal receipt and quota. A force
witness proves the activated instruction and retained anchor instead.
\end{definition}
A proposal, QC, submitted Force, or status map alone does not grant entry.
At Ready, relay the complete bundle to all successor members. Their
complete validation is their own Ready event. Duplicate witnesses do
not reset state. Having entered an old discarded branch does not block
entry to a new session with the same numeric session index.

\begin{protocol}{alg:bootstrap}{Bootstrap the selected region}
Replica $v\in V_{i_{r+1}}$ runs these steps for one validated
session $\sigma_{r+1}$.
Here $\Delta=\Delta_{i_{r+1}}$.
\begin{enumerate}[label=\textbf{\arabic*)},itemsep=3pt]
\item \textbf{Enter.} At Ready, install $A_{r+1,0,0}$,
$\tagg{AUTH}_{r+1,0}$, $\TB_{r+1}$, and $\lambda^{\rm
open}_{r+1}=\lambda_s$. Set the entry lock
to $Q^0_{r+1}$. Enter $(r+1,0,0)$ unfrozen, record its entry time once,
and start progress checks. Send STATUS with this lock and boundary
witness to the initial leader. Process buffered proofs for this session.
\item \textbf{Choose the anchor.} The initial leader waits $2\Delta$
after its own entry. If the view remains current and unfrozen, it sends
NEWVIEW choosing exactly $A_{r+1,0,0}$, justified by $Q^0_{r+1}$ and
the boundary witness.
\item \textbf{Certify and produce.} Members apply the ANCHOR-vote step
of \cref{alg:recovery}. On the ANCHOR QC, the leader starts
\cref{alg:steady} at DATA height $h(A_{r+1,0,0})+1$ and physical
position 1. Later local failures use the same recovery protocol.
\end{enumerate}
\end{protocol}

Let $e=\min_v e_v(r+1,0,0)$ over correct members $v$, the
first correct Ready time. The complete-bundle relay makes
all correct members enter by $e+\Delta$. With a correct initializer and
no quit, NEWVIEW is sent by $e+3\Delta$, its ANCHOR QC reaches the
leader by $e+5\Delta$, and the first DATA proposal reaches all correct
members by $e+6\Delta$. This meets the first progress check. Transfer
before Ready remains subject to the backbone's eventual-delivery rule.

\section{The L1 Force handler}
\label{sec:force}

\subsection{The committee's internal protected prefix}
\label{sec:protection}
The committee maintains a \defterm{protected checkpoint} $G_{\rm msig}$, starting
at the common trusted checkpoint. It stores the complete validated
prefix and its branch provenance. This is internal state, not a
certificate or an extra finality level exposed to clients.
It records a promise: every later Force $F$ satisfies
$\Del(F)\cap G_{\rm msig}=\emptyset$.

\emph{Registration.} On a complete regional-finality bundle, validate the whole missing path
from a trusted checkpoint under the supplied historical L1 snapshot,
parents first. Register $\mu(B)$ once for each newly validated DATA block.
If a Force removed a block while the bundle was in transit, its historical
evidence can still be validated, but the block is only archived. It is
never protected. Ancestors in that path are registered no
later than descendants. A duplicate or a different receipt encoding
does not restart a deadline. The monitoring bound gives
\[
 \mu(B)\le s_u(B)+T_{\rm mon}
\]
for every correct observer whose relay validates that block. An earlier
committee observation may make $\mu(B)$ smaller.

\emph{Protection.} At time $\mu(B)+T_{\rm msig}$, if $B$ remains on the current branch,
the committee protects the next continuous prefix through $B$, unless
$B\in\Del(F)$ for its outstanding Force $F$. Already protected blocks
stay protected. Blocks removed by an activated Force stay archived.
Only registered blocks are protected, so every protected block has direct or covering commitment
evidence and the history needed to serve as a retained anchor.

\emph{Event order.} Protection and actual Force submission
are processed by one coordinated state machine, in time order:
\begin{enumerate}[label=(\arabic*)]
\item before a submission at time $s$, complete every protection
whose deadline is strictly before $s$;
\item at $s$, recheck $\Del(F)\cap G_{\rm msig}=\emptyset$, even if
signatures were collected earlier, and submit;
\item then complete the protections whose deadline is $s$.
\end{enumerate}
Background processing delay never extends a deadline.

While $F$ is outstanding, receipts for blocks in $\Del(F)$, including
future descendants of the replaced regular boundary, do not enter
$G_{\rm msig}$. Receipts for the retained prefix
continue to be processed. After activation, removed blocks remain
archived and cannot regain protection through a status update.

\subsection{A Force instruction and its three stages}
\label{sec:force-stages}

\begin{protocol}{alg:force}{Submit, finalize, and activate a Force}
The committee runs step 1. Every node runs steps 2 and 3.
\begin{enumerate}[label=\textbf{\arabic*)},itemsep=3pt]
\item \textbf{Submit (committee).} On a decision to replace an unavailable region:
\begin{itemize}
\item Finish all strictly overdue protections. Require no outstanding
Force and no unfinalized update. Identify the exact target
session on the current branch. Reject a target that was already
the target of a Force.
\item Choose an available DATA anchor on the target's current
branch, with direct or covering commitment evidence and
$h(A_{r,0,0})\le h_*\le H_r$. Verify
its regional history on the current branch. Require
$\Del(F)\cap G_{\rm msig}=\emptyset$.
\item If protected output already lies beyond the target's regular
boundary, no $h_*$ meets this requirement, and that boundary cannot be replaced. A stale failure report may
mark the old region off for future selection; it cannot remove that
protected output.
\item Keep the anchor, commitment evidence, and needed history available outside
the failed region. A known suitable anchor is enough; no proof of a
globally highest committed block is required.
\item Choose a finite activation height as in \cref{sec:force-stages}.
Mark the target region off in the future status map. Atomically
submit the authorized transaction and record $\Del(F)$. Then
process protections due at the same time.
\end{itemize}
If a Force is already outstanding, retain new failure reports. Re-evaluate
them on the branch present after activation. If only a passive update
is unfinalized, wait for its finalization.

\item \textbf{L1 finalized (node).} On reading a finalized Force record:
\begin{itemize}
\item Verify its authorization, instruction identity, target session,
L1 proof, and activation fields. Use the common finalized L1 history.
If the target session is not yet known, keep the record and
resume this step when its boundary arrives. Then check the target and
$h_*$ (\cref{sec:l1}); ignore an invalid record.
\item Record $\Del(F)$ as its future deletion. This applies even
when $h_*=H_r$.
\item Prevent global promotion of $\Del(F)$ at once. Cache or fetch
the retained anchor for the coming handoff. Ordinary regional BFT and
regular handoffs continue until activation.
\end{itemize}

\item \textbf{L1 activated (node).} Once the activation L1 block is finalized,
apply the status update and invoke the forced path in \cref{alg:rotation}.
Apply the instruction once. An already applied instruction or its retry
does not cause another fork, reset a timer, or veto the new branch.
Keep its historical cut when checking old-branch evidence.
\end{enumerate}
\end{protocol}

The instruction binds the chain, target session ID and index $r$,
retained DATA height $h_*$, status map, and activation L1 height
$\lambda_a$. Its \defterm{identity} is the hash of these canonical fields, not a
transaction hash or a choice of multisignature encoding. Retries keep
the same identity and effect.

\begin{table}[htbp]\centering\small
\caption{Force stages and their effects.}\label{tab:force-stages}
\begin{tabularx}{\linewidth}{@{}p{.19\linewidth}Y Y@{}}\toprule
Stage & Event & Effect\\\midrule
\defterm{Submitted} & The signed instruction is actually handed to the bounded
L1 submission service. & The committee records $\Del(F)$ as a
pending deletion and stops protecting it. Signing alone is not submission.\\
\defterm{L1 finalized} & The instruction appears in finalized L1 records. & All
observers block global promotion of $\Del(F)$. Production
authority is not yet changed.\\
\defterm{L1 activated} & L1 block $\lambda_a$ is finalized. & Apply the status map
and perform the external fork and rotation.\\\bottomrule
\end{tabularx}
\end{table}

\paragraph{Committee rules for updates.}
\begin{itemize}
\item Only one Force may be outstanding from Submitted through activation.
The committee keeps its contents fixed and resubmits as needed. It does
not use a cancellable or replaceable pending Force to postpone protection.
\item The committee submits an update only after every update
it submitted earlier is finalized. L1 order therefore equals submission
order, and a passive update cannot overtake a pending Force. Passive
updates may still be submitted after an outstanding Force is finalized;
they activate after it.
\item It chooses $\lambda_a$ above the activation height of every
earlier update, and large enough that $\lambda_a\ge\lambda_p+N$ holds
for any inclusion height $\lambda_p$ reached within $T_{\rm L1}$.
\item A session is the target of at most one Force. A new
Force may follow after activation, against any session on the
resulting branch that has not been a target, including the session that
the activated Force opened. A later live status does not restore
a deleted suffix.
\end{itemize}
This restriction serializes exceptional Force instructions. Regular
handoffs still run off L1.

\subsection{Processing rules}
\Cref{alg:force} gives the processing rules. A node that first sees the instruction after its activation height
processes finalization and activation in order during that same update.
It does not wait for another message. Malformed or unauthorized records
are rejected (\cref{sec:l1}). An archived
instruction is interpreted only in its original branch context.

The L1 record does not certify the committee's private receipt times.
The committee's protection rule is part of the honest-service assumption.
Observers rely on that rule and the stated time bounds; they need neither
a committee acknowledgement nor a view of its mempool submissions.

\section{The external observer}
\label{sec:observers}
Each observer keeps two prefixes: its regional output
(\cref{sec:regional-finality}) and an
irreversible global prefix ending at $G_u$. It also keeps current
finalized L1 proofs, boundary history, and first relay times.

\begin{protocol}{alg:external}{Verify output and advance global finality}
Observer $u$ runs these steps. At equal time, step 1 precedes step 3.
\begin{enumerate}[label=\textbf{\arabic*)},itemsep=3pt]
\item \textbf{Update L1.} Read and validate the current finalized
records. Run steps 2 and 3 of \cref{alg:force}, including any activation already due.
Keep each pending Force's $\Del(F)$. If the finalized
snapshot is not available, defer global promotion.
Before step~3 runs at time $z$, run this step at $z$, so the
snapshot holds every Force available by $z$ (\cref{ass:monitor}).

\item \textbf{Accept and relay regional output.} On new block evidence:
\begin{itemize}
\item Recover the missing path from $G_u$. Verify
regional finality, including the derived region and session.
Process DATA ancestors before children. Request missing data or receipts.
\item Append valid blocks on the current branch to regional output.
For each block not previously relayed, set $s_u(B)$ when full validation
and relay occur, and send its complete bundle to the committee.
\item Batch or deduplicate transmission without delaying it beyond
the monitoring bound. Keep the original time on duplicates.
\end{itemize}

\item \textbf{Promote the next DATA block.} Starting immediately after
$G_u$, apply \cref{def:global-finality} using the snapshot of step~1.
If the next block lacks evidence, has not waited $\Twait$, or is
in $\Del(F)$ for an applicable finalized Force, stop promotion there.
Otherwise append it to the global prefix and continue in order.

\item \textbf{Maintain the prefixes.} On a Force activation, remove
$\Del(F)$ from the regional output. Archive its evidence. A valid Force preserves $G_u$. Treat any
purported removal of $G_u$ as a violation of the assumed control
service; do not roll back the global prefix.
\end{enumerate}
\end{protocol}

The key for a timer or block record is its session ID and block
hash, not its numeric height alone. New blocks on a forced branch start
new waits. The same retained block keeps its earlier wait. Persistence
prevents replay from restarting waits; a node that cannot recover its
local timing state conservatively waits again.

\subsection{Requests and retained data}
\label{sec:requests}
A request is \defterm{admissible} when it passes the application's admission
checks. Each node queues it in order of first receipt. Repeated delivery
does not change its queue position. New arrivals do not move
ahead of queued requests, and different nodes need not have the same
queue order. A correct leader skips requests
already in its proposed ancestry, including the inherited regional
prefix, and fills the batch with the oldest remaining requests that fit.
If no request can be included, it may propose an empty DATA block.

A request is finally removed from the pending queue only when included
in the global prefix. A request in a discarded regional suffix remains
pending and can be proposed on the new branch. Regional commitment or
regular handoff alone does not remove it. Keep the block data and
receipts needed for recovery and checkpoint verification; global
promotion does not authorize deleting their only available copy.

\section{Examples and guarantees}
\label{sec:examples}

\subsection{Regular rotation and a late DRAIN}
Let $h(A_{r,0,0})=100$, $\TB_r=10$, and $K=4$. The terms produce
DATA 101--104, 105--108, and 109--110. A terminal receipt for DATA 110
permits normal rotation to the selected region, which starts with DATA
111. No global wait is inserted into this handoff.

\Cref{fig:regular-cross-region} traces the communication in
\cref{alg:rotation,alg:bootstrap}. Region A serves session $r$, and the
boundary selects region B for session $r+1$. The old region first
collects FINAL messages into $\FC_r(T_r)$; a proposal QC alone does not
permit this handoff. The complete REGULAR witness crosses the backbone.
Each successor member enters only at its own Ready event, after checking
the witness, application history, and finalized authority. The initial
leader then waits $2\Delta_{\rm B}$, certifies the inherited anchor, and
proposes DATA 111. Global confirmation of DATA 110 runs independently.

\begin{figure}[htbp]\centering
% Protocols 3 and 4: a terminal receipt transfers the turn from A to B.
% Representative arrows only; time runs downward and is not to scale.
\begin{tikzpicture}[
  x=1cm,y=-1cm,font=\fontsize{9}{10.5}\selectfont,
  party/.style={draw,rounded corners=2pt,fill=white,align=center,
    inner sep=3pt,minimum width=2.2cm},
  lane/.style={draw=black!30,densely dashed},
  msg/.style={-{Stealth[length=1.7mm]},semithick,draw=rrBlue},
  backbone/.style={msg,densely dashed},
  tag/.style={fill=white,inner sep=1.5pt,font=\fontsize{8}{9}\selectfont},
  note/.style={anchor=north west,align=flush left,text width=5.55cm,
    inner sep=0pt,font=\fontsize{8.5}{10}\selectfont},
  phase/.style={anchor=west,fill=white,inner sep=2pt,
    font=\fontsize{9}{10.5}\selectfont\bfseries},
  dot/.style={circle,fill=black,inner sep=1.1pt},
  timer/.style={line width=2.6pt,draw=rrTeal!75},
  wait/.style={line width=2.6pt,draw=rrBlue!70!black,
    dash pattern=on 2.5pt off 1.5pt},
  entry/.style={draw=rrBlue,line width=1.2pt,dotted}
]
  \def\xA{1.2}\def\xR{3.8}\def\xL{6.4}\def\xB{9.0}\def\xn{10.0}
  \path[use as bounding box] (-.1,-.65) rectangle (15.6,15.75);
  \node[party] at (\xA,0) {A: leader /\\receipt holder};
  \node[party] at (\xR,0) {A: other\\replicas};
  \node[party] at (\xL,0) {B: initial leader\\$L(r+1,0,0)$};
  \node[party] at (\xB,0) {B: other\\replicas};
  \foreach \x in {\xA,\xR,\xL,\xB}
    \draw[lane] (\x,.55)--(\x,15.35);

  \node[phase] at (0,.85) {1.\ A commits its quota endpoint};
  \draw[msg] (\xA,1.40)--node[tag,above,sloped] {$T_r=110$}(\xR,1.70);
  \draw[msg] (\xR,2.15)--node[tag,above,sloped] {DATA votes}(\xA,2.45);
  \foreach \x/\s/\e in {\xA/1.40/3.80,\xR/1.70/4.10} {
    \draw[timer] (\x+.20,\s)--(\x+.20,\e);
    \node[dot] at (\x,\s) {};
    \draw[thick] (\x-.12,\e)--(\x+.45,\e);
  }
  \node[note] at (\xn,1.35) {Replicas validate and vote for the last DATA
    block. Each starts its own $2\Delta_{\rm A}$ commit timer.};
  \node[tag,anchor=east] at (\xA-.15,2.90) {QC$(110)$};
  \node[note] at (\xn,2.80) {The terminal QC does not hand the turn to B.
    The quota permits no DATA 111 in A's session.};
  \draw[msg] (\xR,4.25)--node[tag,above,sloped] {FINAL$(110)$}(\xA,4.55);
  \node[dot] at (\xA,4.55) {};
  \node[note] at (\xn,3.95) {Timers expire; replicas commit 110 and send
    FINAL. Matching signatures of weight $>F_{\rm A}$ form
    $\FC_r(T_r)$. Receipt holders complete the old engine.};

  \node[phase] at (0,5.55) {2.\ The regular boundary crosses the backbone};
  \node[note] at (\xn,5.50) {Verify the authorized session and quota.
    Derive B from $\chi(\lambda_s)$ and form the REGULAR witness;
    no new L1 transaction is needed.};
  \draw[backbone] (\xA,6.35)--node[tag,pos=.55,above,sloped]
    {REGULAR witness + bundle}(\xL,7.35);
  \draw[backbone] (\xA,6.35)--(\xB,7.55);
  \node[note] at (\xn,6.90) {Transfer has no fixed delay bound. Global
    confirmation of 110 proceeds separately; this handoff does not wait
    for it.};

  \node[phase] at (0,8.05) {3.\ B validates, certifies the anchor, and produces};
  \draw[entry] (\xL-.6,8.45)--(\xL,8.45)--(\xB,8.65)--(\xB+.6,8.65);
  \foreach \x/\y in {\xL/8.45,\xB/8.65} \node[dot] at (\x,\y) {};
  \draw[msg] (\xL,8.75)--node[tag,above,sloped] {Ready relay}(\xB,8.95);
  \draw[msg] (\xB,9.25)--node[tag,above,sloped] {STATUS}(\xL,9.50);
  \draw[wait] (\xL+.20,8.45)--(\xL+.20,10.85);
  \node[note] at (\xn,8.40) {Ready requires the boundary, anchor 110,
    history, and finalized authority. Enter $(r+1,0,0)$; relay the bundle
    and send STATUS.};
  \node[note] at (\xn,9.95) {The initial leader waits $2\Delta_{\rm B}$
    from its own entry. Progress checks already run.};
  \draw[msg] (\xL,10.90)--node[tag,above,sloped] {NEWVIEW$(110)$}(\xB,11.20);
  \draw[msg] (\xB,11.65)--node[tag,above,sloped] {ANCHOR votes}(\xL,11.95);
  \node[dot] at (\xL,11.95) {};
  \node[note] at (\xn,11.05) {NEWVIEW must choose the inherited anchor.
    Weight $>F_{\rm B}$ gives an ANCHOR QC; these votes start no timer.};
  \draw[msg] (\xL,12.55)--node[tag,above,sloped] {DATA 111}(\xB,12.85);
  \foreach \x/\s/\e in {\xL/12.55/14.95,\xB/12.85/15.25} {
    \draw[timer] (\x+.20,\s)--(\x+.20,\e);
    \node[dot] at (\x,\s) {};
    \draw[thick] (\x-.12,\e)--(\x+.45,\e);
  }
  \draw[msg] (\xB,13.35)--node[tag,above,sloped] {DATA votes}(\xL,13.65);
  \node[note] at (\xn,12.55) {The first fresh block has DATA height 111
    and physical position 1. Its votes start new commit timers. A QC
    permits the next proposal before 111 commits.};
  \node[note] at (\xn,14.60) {If uninterrupted, the new timers commit 111;
    replicas send FINAL for B's session.};
\end{tikzpicture}

\caption{Regular cross-region rotation (\cref{alg:rotation,alg:bootstrap}).
Time runs downward; the lanes show representative participants, and
vertical spacing is schematic. FINAL and ANCHOR votes are broadcast
within their respective regions; only representative arrows are drawn.
A terminal receipt opens the regular boundary. Backbone transfer has no
fixed delay bound. Ready starts the new session's entry clock, followed
by NEWVIEW, ANCHOR certification, and the first DATA proposal. Neither
the witness nor the new region waits for global finality.}
\label{fig:regular-cross-region}
\end{figure}
\FloatBarrier

Now suppose DATA 110 has a QC but its timer is cancelled by Quit.
Recovery opens at 110 and produces DRAIN. When DRAIN commits 110,
replicas send FINAL for 110. Its old L1 reference stays unchanged.
An observer that first validates and relays that receipt at time 100
starts $s_u(110)=100$, even if the block's L1 timestamp was 0. It waits
until $100+\Twait$. A Force submitted at 95 and finalized at 105 is seen
during that wait and prevents global promotion if it cuts below 110.
There is no need to protect an uncommitted QC from time 0.

\subsection{A Force submitted before the committee deadline}
Suppose a correct observer validates and relays DATA 107 at time 100.
Let $T_{\rm mon}=2$, $T_{\rm msig}=5$, and $T_{\rm L1}=3$, so
$\Twait=10$. The committee validates it by 102 and decides by 107.
A Force submitted at 107 is finalized by 110. At 110 the observer
processes that record before its global test. If the Force keeps only
DATA 106, it cannot promote 107, even if activation is at 115. At
activation, 107 and its old descendants are removed. The new region may
produce a different 107 under a new session, with a new observer
wait.

If no Force is submitted by the committee's actual deadline, its
protected prefix advances through 107. A later Force must preserve it.
The observer can promote 107 at 110 if its predecessors are already
global and the current L1 history passes the checks. The committee sends
no acknowledgement in either case.

\begin{figure}[htbp]\centering
\begin{tikzpicture}[x=1cm,y=1cm,font=\small,
  box/.style={draw,rounded corners=2pt,align=center,text width=3.65cm,
  minimum height=1.2cm,inner sep=4pt},
  arr/.style={-{Stealth[length=2mm]},semithick}]
\node[box,fill=rrBlue!7] (s) at (0,0) {Submitted\\committee records the cut};
\node[box,fill=rrRed!7] (f) at (5.25,0) {L1 finalized\\affected blocks cannot be global};
\node[box,fill=rrTeal!8] (a) at (10.5,0) {L1 activated\\external branch switches};
\draw[arr] (s)--node[above,yshift=8mm,font=\scriptsize]{at most $T_{\rm L1}$}(f);
\draw[arr] (f)--node[above,yshift=8mm,font=\scriptsize,align=center]{activation block\\finalizes}(a);
\end{tikzpicture}
\caption{Three stages of one Force. The global waiting period covers
monitoring, the decision window, and Force finalization. A finalized
instruction blocks promotion before the external fork activates.}
\label{fig:force-stages}
\end{figure}

\paragraph{Example: forced cross-region handoff.}
\Cref{fig:forced-cross-region} expands the three stages into the
communication of \cref{alg:force,alg:rotation,alg:bootstrap}. The
committee targets region A's session and retains DATA 106. Before
submission it checks that the entire deletion scope, including regular
successor sessions, avoids its protected prefix. The Force binds the
target session ID, cut, future status map, and activation height. Once
the record is finalized, nodes block global promotion of its deletion
scope but keep following the old branch until activation. At activation
they remove that scope, obtain the retained anchor and history, and
select region B. B enters only after validating the FORCE witness and
its authority; it then uses the same bootstrap as a regular successor
and proposes a new DATA 107. The old regional engines keep their own
state and timers.

\begin{figure}[htbp]\centering
% Protocols 5, 3, and 4: Force timeline followed by successor bootstrap.
% Dashed arrows are finalized L1 reads or backbone transfer, as labeled.
\begin{tikzpicture}[
  x=1cm,y=-1cm,font=\fontsize{8.5}{10}\selectfont,
  party/.style={draw,rounded corners=2pt,fill=white,align=center,
    inner sep=3pt,text width=1.55cm,minimum height=.65cm},
  lane/.style={draw=black!30,densely dashed},
  msg/.style={-{Stealth[length=1.7mm]},semithick,draw=rrBlue},
  read/.style={msg,densely dashed},
  force/.style={msg,draw=rrRed},
  tag/.style={fill=white,inner sep=1.5pt,font=\fontsize{8}{9}\selectfont},
  note/.style={anchor=north west,align=flush left,text width=5.55cm,
    inner sep=0pt,font=\fontsize{8.5}{10}\selectfont},
  phase/.style={anchor=west,fill=white,inner sep=2pt,
    font=\fontsize{9}{10.5}\selectfont\bfseries},
  dot/.style={circle,fill=black,inner sep=1.1pt},
  pending/.style={line width=2.6pt,draw=rrRed!60},
  wait/.style={line width=2.6pt,draw=rrBlue!70!black,
    dash pattern=on 2.5pt off 1.5pt}
]
  \def\xC{1.0}\def\xI{3.0}\def\xN{5.0}\def\xL{7.0}\def\xB{9.0}\def\xn{10.0}
  \path[use as bounding box] (-.1,-.65) rectangle (15.6,17.40);
  \node[party] at (\xC,0) {Multisig\\committee};
  \node[party] at (\xI,0) {L1};
  \node[party] at (\xN,0) {Node /\\boundary relay};
  \node[party] at (\xL,0) {B: initial\\leader};
  \node[party] at (\xB,0) {B: other\\replicas};
  \foreach \x in {\xC,\xI,\xN,\xL,\xB}
    \draw[lane] (\x,.55)--(\x,16.85);

  \node[phase] at (0,.90) {1.\ Submit the Force; finalize its record};
  \node[dot] at (\xC,1.55) {};
  \node[note] at (\xn,1.35) {Validate anchor 106 and submission rules.
    Require $\Del(F)\cap G_{\rm msig}=\emptyset$. Keep anchor data
    available outside A.};
  \draw[force] (\xC,2.65)--node[tag,above,sloped] {Force call}(\xI,2.95);
  \draw[pending] (\xC+.20,2.65)--(\xC+.20,7.70);
  \node[note] at (\xn,2.75) {Submission at $s$ records the fixed deletion
    scope and stops its protection. Signing alone is not submission.};
  \node[dot] at (\xI,3.95) {};
  \node[tag,anchor=east] at (\xI-.15,3.95) {included at $\lambda_p$};
  \node[note] at (\xn,4.00) {The record binds target ID, $h_*=106$,
    future status map, and $\lambda_a\ge\lambda_p+N$.};
  \draw[read,draw=rrRed] (\xI,5.00)--node[tag,pos=.50,above,sloped]
    {finalized Force}(\xN,5.30);
  \node[dot] at (\xN,5.30) {};
  \node[note] at (\xn,5.00) {By $s+T_{\rm L1}$, correct nodes can read
    the finalized record. Validate it and block global promotion of
    $\Del(F)$ immediately.};
  \node[tag,align=center] at (5.0,6.45)
    {Until activation:\\old branch continues};
  \node[note] at (\xn,6.35) {Regional BFT and regular handoffs still run.
    Nodes may prepare; production authority has not switched.};

  \node[phase] at (0,7.30) {2.\ Activate the cut; construct the forced boundary};
  \draw[read,draw=rrRed] (\xI,7.75)--node[tag,pos=.52,above,sloped]
    {activate at $\lambda_a$}(\xN,8.05);
  \node[dot] at (\xN,8.05) {};
  \node[note] at (\xn,7.75) {Remove A's DATA above 106 and its regular
    successor sessions. Apply the status map; preserve $G_u$ and archive
    evidence. Old engine timers are unchanged.};
  \draw[msg] (\xN,9.10)--node[tag,pos=.48,above,sloped]
    {fetch anchor / history}(\xC,9.40);
  \draw[msg] (\xC,9.75)--node[tag,pos=.52,above,sloped]
    {106 + commitment evidence}(\xN,10.05);
  \node[note] at (\xn,9.50) {Fetch retained data; the committee is one
    possible source. Select B at the least eligible height at or above
    $\max\{\lambda_a,\lambda(106)\}$.};
  \draw[read] (\xN,10.85)--node[tag,pos=.50,above,sloped]
    {FORCE witness}(\xL,11.25);
  \draw[read] (\xN,10.85)--(\xB,11.45);
  \node[note] at (\xn,11.00) {Send the complete boundary bundle over
    the backbone. Missing data, a live successor, or finalized authority
    can delay Ready.};

  \node[phase] at (0,12.00) {3.\ B starts a new session from DATA 106};
  \foreach \x/\y in {\xL/12.50,\xB/12.65} \node[dot] at (\x,\y) {};
  \draw[msg] (\xL,12.70)--node[tag,above,sloped] {Ready relay}(\xB,12.90);
  \draw[msg] (\xB,13.25)--node[tag,above,sloped] {STATUS}(\xL,13.50);
  \draw[wait] (\xL+.20,12.50)--(\xL+.20,14.20);
  \node[note] at (\xn,12.50) {At Ready, install anchor 106, the new ID,
    quota, and authority. Enter $(r+1,0,0)$; send STATUS. The leader waits
    $2\Delta_{\rm B}$ from its own entry.};
  \draw[msg] (\xL,14.25)--node[tag,above,sloped] {NEWVIEW$(106)$}(\xB,14.55);
  \draw[msg] (\xB,15.00)--node[tag,above,sloped] {ANCHOR votes}(\xL,15.30);
  \node[dot] at (\xL,15.30) {};
  \node[note] at (\xn,14.45) {Certify the exact inherited anchor with
    weight $>F_{\rm B}$. ANCHOR votes start no commit timer.};
  \draw[msg] (\xL,16.05)--node[tag,above,sloped] {new DATA 107}(\xB,16.35);
  \node[note] at (\xn,16.00) {The ANCHOR QC permits a fresh 107. DATA
    votes start timers. Validation and relay of its commitment evidence
    start new observer waits.};
\end{tikzpicture}

\caption{Forced cross-region rotation. The submission-to-finalization
bound $T_{\rm L1}$ ends at the finalized Force record, not at activation.
Dashed arrows show finalized L1 reads and backbone transfer. Finalization
vetoes global promotion; finalization of L1 block $\lambda_a$ switches
the external branch. Anchor retrieval, successor selection, backbone
transfer, and authority availability can delay Ready. After Ready, the
new region performs STATUS, NEWVIEW, ANCHOR voting, and DATA proposal
steps. Time runs downward and is not to scale; vote broadcasts are
represented by selected arrows.}
\label{fig:forced-cross-region}
\end{figure}
\FloatBarrier

\clearpage
\subsection{An external observer's finality timeline}
\Cref{fig:observer-finality} traces \cref{alg:external} for the same
DATA 107. Set $s_u(107)=100$ and $\mu(107)=102$, giving protection
deadline 107 and global test 110. Panel (a) follows a Force submitted at
that deadline and activated at time 115. Panel (b) follows a separate
execution without a Force covering the block by its deadline.

\begin{figure}[htbp]\centering
% Protocol 6 and the committee's protection rule. Two distinct executions.
% Numerical example: s_u=100, mu=102, decision deadline=107, test=110.
\begin{tikzpicture}[
  x=1cm,y=-1cm,font=\fontsize{9}{10.5}\selectfont,
  party/.style={draw,rounded corners=2pt,fill=white,align=center,
    inner sep=3pt,minimum width=2.15cm},
  lane/.style={draw=black!30,densely dashed},
  msg/.style={-{Stealth[length=1.7mm]},semithick,draw=rrBlue},
  read/.style={msg,densely dashed,draw=rrRed},
  tag/.style={fill=white,inner sep=1.5pt,font=\fontsize{8}{9}\selectfont},
  note/.style={anchor=north west,align=flush left,text width=5.55cm,
    inner sep=0pt,font=\fontsize{8.5}{10}\selectfont},
  phase/.style={anchor=west,fill=white,inner sep=2pt,
    font=\fontsize{9}{10.5}\selectfont\bfseries},
  dot/.style={circle,fill=black,inner sep=1.1pt},
  wait/.style={line width=2.6pt,draw=rrTeal!75},
  decision/.style={line width=2.6pt,draw=rrBlue!65},
  outcome/.style={draw=rrTeal,rounded corners=2pt,fill=rrTeal!5,
    align=center,text width=4.1cm,minimum height=1.9cm,inner sep=4pt,
    font=\fontsize{8.5}{10}\selectfont}
]
  \def\xA{1.2}\def\xU{3.8}\def\xC{6.4}\def\xI{9.0}\def\xn{10.0}
  \path[use as bounding box] (-.1,-.65) rectangle (15.6,16.45);
  \node[party] at (\xA,0) {Regional evidence\\source};
  \node[party] at (\xU,0) {External\\observer $u$};
  \node[party] at (\xC,0) {Multisig\\committee};
  \node[party] at (\xI,0) {L1};
  \foreach \x in {\xA,\xU,\xC,\xI}
    \draw[lane] (\x,.55)--(\x,12.35);

  \node[phase] at (0,.85) {(a)\ A Force cuts DATA 107 before it can become global};
  \draw[msg] (\xA,1.30)--node[tag,above,sloped]
    {107 + evidence}(\xU,1.75);
  \node[dot] at (\xU,1.75) {};
  \draw[wait] (\xU+.20,1.75)--(\xU+.20,9.80);
  \draw[msg] (\xU,1.75)--node[tag,pos=.60,above,sloped]
    {backbone relayed}(\xC,3.05);
  \node[tag,anchor=east] at (\xU-.15,2.25) {$s_u(107)=100$};
  \node[note] at (\xn,1.40) {Validate 107 on the current branch and relay
    its complete evidence. Start the $\Twait=10$ wait. Old L1 timestamps
    do not start it.};
  \node[dot] at (\xC,3.05) {};
  \draw[decision] (\xC+.20,3.05)--(\xC+.20,6.60);
  \node[tag,anchor=east] at (\xC-.15,3.55) {$\mu(107)=102$};
  \node[note] at (\xn,3.00) {Monitoring takes at most $T_{\rm mon}=2$.
    Here validation is at 102, so the committee's decision window
    $T_{\rm msig}=5$ ends at 107.};
  \node[tag,align=center] at (\xU,5.05)
    {Regional production /\\regular handoff may continue};
  \node[note] at (\xn,4.80) {The observer wait and the committee's
    decision window overlap production. No committee acknowledgement
    is required.};
  \draw[msg,draw=rrRed] (\xC,6.60)--node[tag,above,sloped]
    {Force: retain 106}(\xI,6.90);
  \node[dot] at (\xC,6.60) {};
  \node[tag,anchor=east] at (\xC-.15,7.10) {$d=107$};
  \node[note] at (\xn,6.55) {At the actual deadline, submission precedes
    protection. The outstanding Force excludes 107 and its descendants
    from the protected prefix.};
  \node[note] at (\xn,8.05) {$T_{\rm L1}=3$: the submitted Force is
    finalized and available by 110. This bound does not include
    activation.};
  \draw[read] (\xI,9.20)--node[tag,pos=.48,above,sloped]
    {finalized Force proof}(\xU,9.80);
  \node[dot,fill=rrRed] at (\xU,9.80) {};
  \node[tag,anchor=east,text=rrRed] at (\xU-.15,10.30)
    {$z=110$: no promotion};
  \node[note] at (\xn,9.45) {P6 step 1 reads L1 before step 3. Although
    the wait has ended, $107\in\Del(F)$, so global promotion stops.
    The regional branch has not yet switched.};
  \draw[read] (\xI,11.35)--node[tag,pos=.48,above,sloped]
    {activation block finalized}(\xU,11.85);
  \node[tag,anchor=east] at (\xU-.15,12.15) {activation at 115};
  \node[note] at (\xn,11.55) {P6 step 4 removes the old regional suffix
    and archives its evidence. $G_u$ is preserved. A new 107 on the
    forced branch has a new validation-and-relay wait.};

  \node[phase] at (0,13.05) {(b)\ Separate execution: no Force covering 107 by its deadline};
  \node[outcome] (relay) at (2.45,14.55)
    {$100$: observer validates and relays 107.\\
     Its waiting period starts.};
  \node[outcome] (protect) at (7.75,14.55)
    {$107$: committee protects the continuous prefix through 107.\\
     Later Forces must preserve it.};
  \node[outcome] (global) at (13.05,14.55)
    {$110$: observer may promote 107.\\
     Predecessors must be global; current L1 and branch checks must pass.};
  \draw[msg,draw=rrTeal] (relay.east)--(protect.west);
  \draw[msg,draw=rrTeal] (protect.east)--(global.west);
  \node[tag,align=center] at (7.75,16.00)
    {$\Twait=T_{\rm mon}+T_{\rm msig}+T_{\rm L1}=2+5+3=10$};
\end{tikzpicture}

\caption{External finality with
$T_{\rm mon}=2$, $T_{\rm msig}=5$, and $T_{\rm L1}=3$.
The observer starts its own $\Twait=10$ wait only upon complete
validation and relay. Panel (a) shows a Force submitted at the committee
deadline and finalized by the global test: L1 processing comes first,
so DATA 107 is not promoted. Activation later removes the regional
suffix while preserving $G_u$. Panel (b) is a separate execution without
a Force covering the block by that deadline; protection precedes the
observer's conditional promotion. Time runs downward in panel (a) and
left to right in panel (b); spacing is schematic.}
\label{fig:observer-finality}
\end{figure}
\FloatBarrier

\subsection{Conditional protection argument}
\label{sec:safety-argument}
Assume regional safety and receipt validity within each session,
the stated honest committee rules, and the monitoring and L1 bounds.
We show by induction over $u$'s promotions, starting from its
trusted checkpoint, that no valid Force removes a block from a correct
observer's global prefix.

Suppose $u$ promotes $B$ at time $z\ge s_u(B)+\Twait$. Let
$d=\mu(B)+T_{\rm msig}$ be the committee's deadline for $B$. Since
$\mu(B)\le s_u(B)+T_{\rm mon}$,
\[
 d+T_{\rm L1}\le s_u(B)+\Twait\le z.
\]
Let $F$ be any Force with $B\in\Del(F)$.
\begin{itemize}
\item \emph{$F$ was submitted at or before $d$.} Then $F$ is
finalized and available by $d+T_{\rm L1}\le z$. Step~1 of
\cref{alg:external} reads it before step~3, so condition~(iv) of
\cref{def:global-finality} fails; if $F$ has already activated, $B$ has
no regional finality. Either way, $u$ does not promote $B$.
\item \emph{$F$ was submitted after $d$.} By the first case, no Force
outstanding at $d$ contains $B$, and $B$, which is regionally final at
$z$, was not removed before $d$. So the committee protected $B$ at $d$.
A Force submitted later satisfies $\Del(F)\cap G_{\rm msig}=\emptyset$,
which contradicts $B\in\Del(F)$.
\end{itemize}
The same relay registers $B$'s ancestors no later than $B$, so the
protected prefix through $B$ is continuous.

Equal-time submission precedes protection, and equal-time L1
finalization precedes observer promotion. There is no gap at either
endpoint. The single outstanding Force has a fixed, irrevocable effect;
an ignored or cancelled placeholder cannot extend a protection deadline.
Induction from the trusted checkpoint therefore prevents an accepted
Force from deleting a correct observer's global prefix. Regional safety
and deterministic rotation supply the common retained branch between
Forces. This argument is conditional on those properties; it is not a
new proof of the weighted regional BFT.

\subsection{Progress conditions and limits}
The local engine and regular handoff retain their original progress
conditions. A global prefix grows when complete regional evidence keeps
arriving, monitoring meets its bound, and a working branch remains
uncut long enough to pass $\Twait$. Observers need not confirm at the same
time. Unbounded backbone delay still prevents a fixed end-to-end bound
from regional commitment to every observer's confirmation.

Force recovery also needs eventual activation, an available retained
anchor, a live successor, and its authority. The bound $T_{\rm L1}$ ends
at Force finalization; it is not a bound on reaching any arbitrarily
chosen future activation height. Only one live region is insufficient
for indefinite quota rotations because NextLive excludes the current
region. Repeated failures can also prevent a stable global extension.
These are progress conditions, not reasons to revoke an existing global
prefix.

\subsection{Open design questions}
The following points need a design decision and are not fixed by this
document.
\begin{enumerate}[label=(\arabic*)]
\item An engine whose session was cut or discarded keeps running until it
completes. When it may stop is not specified.
\item The format of the registry proof in the L1 reference $\rho$.
\item A Byzantine leader may choose an old $\lambda(T_r)$, bounded below
by $\lambda^{\rm open}_r$, to influence successor selection.
\item With only one live region, quota rotation stops, because
$\tagg{NextLive}$ excludes the current region.
\item How L1 publishes $\tagg{AUTH}_{r,g}$ for future session indices
without one L1 transaction per regular handoff.
\item The margin in $\lambda_a$ that the committee needs depends on the
L1 block rate, which the model does not bound.
\end{enumerate}

\begingroup\small
\begin{thebibliography}{9}
\bibitem{part1}
Sync HotStuff RR, Part I: Core Protocol. Design document, 13 October 2026
(\texttt{sync\_hotstuff\_rr\_part1\_core.tex}).
\bibitem{sync}
I. Abraham, D. Malkhi, K. Nayak, L. Ren, and M. Yin.
\emph{Sync HotStuff: Simple and Practical Synchronous State Machine Replication}.
IEEE Symposium on Security and Privacy, 2020, pp. 106--118.
\href{https://conferences.computer.org/sp/pdfs/sp/2020/349700a106.pdf}{Published paper}.
\end{thebibliography}
\endgroup

\end{document}
